\documentclass[11pt,a4paper]{ltjsarticle}
\usepackage[margin=23mm]{geometry}
\usepackage{amsmath,amssymb,amsthm,mathtools,booktabs,mathrsfs}
\usepackage[haranoaji]{luatexja-preset}
\usepackage[unicode,hidelinks]{hyperref}
\newtheorem{theorem}{定理}[section]
\newtheorem{proposition}[theorem]{命題}
\newtheorem{lemma}[theorem]{補題}
\newtheorem{corollary}[theorem]{系}
\theoremstyle{definition}\newtheorem{remark}[theorem]{注意}
\newtheorem{definition}[theorem]{定義}
\renewcommand{\proofname}{証明}
\newcommand{\PP}{\mathbb P}\newcommand{\kk}{\mathbb C}
\newcommand{\Sym}{\operatorname{Sym}}\newcommand{\qgr}{\operatorname{qgr}}
\newcommand{\Ext}{\operatorname{Ext}}\newcommand{\Extop}{\operatorname{Ext}}
\newcommand{\Hom}{\operatorname{Hom}}\newcommand{\HH}{\operatorname{HH}}
\newcommand{\Spec}{\operatorname{Spec}}\newcommand{\PGL}{\operatorname{PGL}}
\newcommand{\OO}{\mathcal O}\newcommand{\rank}{\operatorname{rank}}
\newcommand{\Pic}{\operatorname{Pic}}\newcommand{\Aut}{\operatorname{Aut}}
\newcommand{\End}{\operatorname{End}}\newcommand{\RHom}{\operatorname{RHom}}
\newcommand{\Cw}{\mathfrak C^{\rm w}}\newcommand{\Cu}{\mathfrak C^{\rm u}}
\title{非可換三次曲面の二段階の標識忘却\\\large 反標準 heart・有限 Weyl 群・線束の選択}
\author{植田一石}
\date{2026年11月18日 / unmark-v10}
\begin{document}\maketitle
\begin{abstract}
AOU/OU 型標識から弱標識へ、弱標識から無標識へ至る二段階の忘却を、
反標準 heart を基準として集合論的に研究する。
第一段階は固定した弱標識上の標準標識の有限 Weyl 群 $W(E_6)$ による自由度、
第二段階は固定した反標準 heart 内の線束の選択である。
AS heart は標準表示との比較にのみ用いる。
固定圏内の線束選択と、exact 自己同値による軌道商を区別する。

内在的な公理の下で rank-one torsion-free 対象の自己第二 Ext は消え、
例外性は自己 Euler 形式 $1$ と同値である。
三次曲面型の数値表示と構造対象の数値階数 $1$ を仮定すると、
反標準 Serre 消滅、数値階数と内在的 torsion 部分圏の対応、および
有効次数正値性が従い、各数値的線束類の対象は高々一つである。
AOU 三ブロック型の九線束から格子標識を復元し、固定弱標識上の標識数を
$51840$ 以下に抑える。全 Weyl 作用の実現は一種類の三対象交換へ還元される。
零次数 root 型単純対象が存在する場合には、構造対象と同じ反標準切断環を持つが
豊富でない線束ができるという条件付き障害も示す。その実在は未解決である。
全数値類の存在と全線束の許容性は別に確認すべき問題として残す。
可換 locus と許容形式近傍では通常の有限 Weyl 標識を記述し、
一次 Picard 障害によって線束の持上げと自己同値の持上げを区別する。
一般の AOU/OU 部分での有限 Weyl 軌道の実現・網羅性と、
全数値的線束類の実現は、証明された範囲と未解決の範囲を分けて扱う。
\end{abstract}
{\small\tableofcontents}
\clearpage
\section{はじめに}
本稿の基本問題は
\[
 \text{AOU/OU marked}\longrightarrow
 \text{weakly marked}\longrightarrow\text{unmarked}
\]
という二段階の忘却である。
基準とする Abel 圏は反標準環の $\qgr$ として得る
\emph{anticanonical heart} であり、比較する別の心は \emph{AS heart} のみとする。
対象の同型類と圏の exact 同値類の集合を先に分類し、
群oid・stack および相対的な降下の定式化はその後の課題とする。

\paragraph{第一段階：有限 Weyl 群による標識の選択.}
AOU\cite{AOU}とOU\cite{OU2025}の標準表示から
反標準 heart と構造対象が得られる範囲を取り、その弱標識を固定する。
残る標識の自由度を有限 Weyl 群 $W(E_6)$ で記述するのが第一の目標である。
可換三次曲面の通常標識およびその許容形式持上げは既存の比較で扱える。
非可換全域における幾何学的標識の実現・網羅性は、
数値格子上の Weyl 作用があることだけからは従わない。
AOU/OU 標識の像が任意の弱標識付き nc cubic を尽くすとも仮定しない。

\paragraph{第二段階：固定反標準 heart 内の線束の選択.}
弱標識とは $(\mathcal A_{\rm ac},L)$ という構造対象の選択である。
従って最初に調べる集合は $\operatorname{Line}_{\rm adm}(\mathcal A_{\rm ac})$ である。
対象同士のテンソル積や、Picard 格子による固定圏上の作用を先に仮定しない。
各数値的線束類 $v_D=(1,D,D^2/2)$ についての存在・一意性を証明し、
全線束が弱標識として許容されることを確認する。
exact 自己同値による軌道は、線束の選択集合を決定した後に別に取る。

線束対象の数値的判定（定理\ref{thm:intrinsic-line-criterion}）を出発点とする。
定理\ref{thm:ac-serre-positivity}では、全対象に対する三次曲面型 Riemann--Roch と
$r(\OO_0)=1$ の正規化から、反標準 Serre 消滅と
$\mathcal A_1=\{r=0\}$ を導く。有効次数正値性はその帰結であり、
独立に仮定せずに数値的線束類の一意性
（定理\ref{thm:numerical-line-uniqueness}）を得る。
この一意性を用いて、AOU 型の九対象から復元する格子標識への単射と
標識数の上界 $51840$ を証明する（定理\ref{thm:weyl-frame-injection}）。
六つの単純反射のうち五つは直交ブロック内の置換であり、残る一つは
三対象交換として明示できる（系\ref{cor:weyl-one-move}）。
交換後の標識の存在が分かれば、Coxeter 関係と網羅性は一意性から従う。
許容線束を変えて反標準環を作り直しても、Artin--Zhang 再構成により
同じ反標準 heart を得る（命題\ref{prop:ac-line-reconstruction}）。
これは既知の再構成定理の直接の帰結であり、新しい一般的再構成を主張しない。
一方、命題\ref{prop:root-simple-obstruction}は零次数 root 型単純対象の
存在を仮定して、例外性と切断環の一致だけでは固定 heart 内の豊富性を
判定できないことを示す。この単純対象の実例は構成していない。
以上の導出は本稿に証明を記載するが、一般的な先行研究との重複・新規性は
引き続き照合を要する。

\paragraph{先行研究と保持する補助結果.}
AOU の標準九線束\cite[(4.1), (9.44)]{AOU}と古典的格子標識の
$W(E_6)$-torsor\cite[\S2.6]{DGK}は既知の入力である。
本稿の有限上界は、それらと固定反標準 heart 内の線束の一意性を組み合わせた帰結である。
OU\cite[Definition 1.1(3)]{OU2025}の一般の marking は very strong helix の選択であり、
今回の有限性はその全体ではなく、明示した三ブロック線束型に対して述べる。
AZ 豊富性には、Polishchuk\cite[Lemma 2.1(ii)]{PolishchukProj}で明記される
任意 tail の異なる指数を用いた有限生成を使う。
一般の豊富対だけに高次 Ext 消滅を読み込むことはしない。
Serre 双対性と階数一の非零射の単射性から第二 Ext を消す論法には、
Hoffmann--Reede\cite[Lemma 3.9]{HR}の先行例がある。
同補題は $\PP^2$ 上の terminal del Pezzo order に対し、同じ第一 Chern 類の
rank-one torsion-free 加群間の第二 Ext 消滅を示す。
本稿はその分類結果を移し替えるのではなく、指定された反標準 heart の公理から
必要な階数認識と次数正値性を導いて同様の消滅論法を適用する。
対象範囲の包含や一般的な優先権は主張しない。
可換 del Pezzo 曲面の例外層の $K_0$ 類による一意性も既知であり、
Karpov--Nogin\cite[\S1.3, p.~433]{KN}に明記されている。
また root 型単純対象からの評価核は、既知の dual spherical twist の
評価複体\cite[Definition 2.7]{ST}と同じ形を持つ。
本稿ではその一般理論を証明の入力とせず、固定 heart 内の短完全列から
非生成性と切断環の同型を直接確認する。
さらに、AOUの関係式表示、OUのAS標識、Bondal--Orlovの可換自己同値、
LowenとDe Deken--Lowenの変形理論、TodaのFourier--Mukai変形作用を入力とする。
Picard一次障害、非常に一般の整数核、および一次不変式の証明を保持するが、
それらを線束の存在分類や全域的な有限Weyl商の代わりにはしない。
普遍有限代数、支持によるHochschild計算、数値Serre--$E_6$、
形式変形の比較の証明は付録に置く。
数値的 $E_6$ の一致も、AS表示の導来同値も、それだけで
反標準 heart の一致を与えるわけではない。

\section{反標準 heart を基準とする二段階の忘却}\label{sec:ac-set-framework}
本稿の基準は \emph{anticanonical heart}（反標準 heart）である。
AOU/OU の標準表示を使う場合だけ \emph{AS heart} を併用する。
この二つの指定された heart の比較と、標識の忘却とを混同しない。
本稿の分類対象は、まず対象の同型類および圏の exact 同値類の\emph{集合}である。

\subsection{反標準環と固定する heart}
Serre 関手 $\mathsf S$ を持つ導来圏と対象 $e$ に対し $H=\mathsf S^{-1}[2]$ と置く。
反標準切断環と反標準 heart を
\begin{equation}\label{eq:ac-ring-heart}
 R_e=\bigoplus_{n\ge0}\Hom(e,H^ne),\qquad
 \mathcal A_{\rm ac}=\qgr R_e
\end{equation}
と書く。積は $f:e\to H^me$, $g:e\to H^ne$ に対する
$H^m(g)\circ f$ から定める（左右加群の規約を一貫して固定する）。
以下はこの $\qgr$ 構成が noetherian な Abel 圏を与え、
必要な導来比較を使うときにはその指定された実現関手がある範囲で行う。
Hilbert 級数だけからその範囲を拡張しない。
AS heart は AOU/OU の AS-regular $\mathbb Z$-代数 $A^{\rm AS}$ に対する
$\mathcal A_{\rm AS}=\qgr A^{\rm AS}$ を意味する。
$R_e$ と $A^{\rm AS}$、および二つの heart は、名前だけで同一視しない。

先に Abel 圏 $\mathcal A$ と許容構造対象 $e$ が与えられ、
$(e,H)$ が Artin--Zhang 豊富である場合には、
その再構成定理により $\mathcal A\simeq\qgr R_e$ となる
\cite[Theorem 4.5]{AZ}。
従って本稿の内在的公理（定義\ref{def:weak-cubic}）は、
このように反標準 heart として認識された $\mathcal A$ 上で用いる。
単に $H$ が $t$-exact というだけで、任意に選んだ heart を
反標準 heart と呼び替えることはしない。

\subsection{固定した心の中の二つの選択集合}
$\mathcal A=\mathcal A_{\rm ac}$ を固定する。線束対象を
rank-one torsion-free exceptional 対象と呼び、
\[
 \operatorname{Line}(\mathcal A)=
 \{L\in\mathcal A:L\text{ は線束対象}\}/\cong
\]
とする。さらに $(L,H)$ が豊富で所定の切断数を持つものを
$\operatorname{Line}_{\rm adm}(\mathcal A)$ とする。
weak marking とはその一つを選ぶこと、すなわち $(\mathcal A,L)$ である。
全線束の許容性が証明された範囲では二集合は一致する。
一致をまだ証明していない範囲では、その未解決点を明記する。

許容線束 $L$ に対して、$L$ に正規化した AOU/OU 型標識の選択集合を
$\operatorname{Mark}_{\rm AOU/OU}(\mathcal A,L)$ と書く。
ここで、指定した構造対象が $L$ に対応する標準 helix/thread の実現を選び、
固定圏内の対象の同型で同一視する。矢の基底は標識に含めない。
AS 表示から得る場合は、その対象を反標準 heart に移せることを先に確認する。
全域で標準 helix が存在するとは仮定しない。

このとき二段階の分類問題は
\begin{equation}\label{eq:set-two-forgettings}
 (\mathcal A,L,\mu)\longmapsto(\mathcal A,L)\longmapsto\mathcal A,
 \qquad \mu\in\operatorname{Mark}_{\rm AOU/OU}(\mathcal A,L)
\end{equation}
である。固定した対象上では
\begin{equation}\label{eq:set-union-markings}
 \operatorname{Mark}_{\rm AOU/OU}(\mathcal A)
 =\coprod_{[L]\in\operatorname{Line}_{\rm adm}(\mathcal A)}
   \operatorname{Mark}_{\rm AOU/OU}(\mathcal A,L).
\end{equation}
これは選択集合の記述であり、右辺の各項の非空性を主張するものではない。

第一段階の基本問題は、固定した $(\mathcal A,L)$ 上の標識を
\emph{有限 Weyl 群 $W(E_6)$} の作用で記述することである。
第二段階の基本問題は、固定した $\mathcal A_{\rm ac}$ の
\emph{線束の選択}を分類することである。
有限 Weyl 群を線束の全選択を表す群として用いない。
また全線束の同型類が格子で添字付けられることと、
その格子が固定圏の自己同値として作用することは別である。

\begin{proposition}[許容線束による反標準表示の変更]\label{prop:ac-line-reconstruction}
固定した反標準 heart $\mathcal A$ の許容線束 $L$ に対して
\[
 R_L=\bigoplus_{n\ge0}\Hom_{\mathcal A}(L,H^nL)
\]
と置く。Artin--Zhang 再構成の仮定の下で
\[
 \mathcal A\xrightarrow{\sim}\qgr R_L
\]
という exact 同値があり、$L$ は構造対象 $\pi R_L$ に対応する。
従って weak marking の変更は、固定した heart の中で
構造対象を変える操作として扱える。
これは $R_L$ と $R_{L'}$ の次数付き環同型も、
$L$ と $L'$ を交換する $\mathcal A$ の自己同値も主張しない。
\end{proposition}
\begin{proof}
$L$ は noetherian で、Hom は有限次元、$(L,H)$ は豊富である。
Artin--Zhang 再構成\cite[Theorem 4.5]{AZ}をこの組に適用する。
各 $L$ で同じ $\mathcal A$ を再構成しているので、結論を得る。
二つの構造対象を交換する自己同値は、この再構成とは別の条件である。
\end{proof}

\subsection{自己同値で割る操作は後に分けて行う}
上では固定圏内の線束を、対象の同型だけで同一視した。
圏の同定まで忘れて weakly marked 対象の exact 同値類を取るときには、
さらに
\begin{equation}\label{eq:set-exact-orbits}
 \operatorname{Line}_{\rm adm}(\mathcal A)/\operatorname{Aut}_{\rm ex}(\mathcal A)
\end{equation}
という\emph{集合としての軌道商}を考える。
ここでも最初から線束の選択集合を軌道商で置き換えない。
まず各数値類での存在・一意性と許容性を解き、次にこの作用を計算する。
群oid・stack・相対降下は、この集合論的分類の後の課題とする。

% !TeX root = ../nc-cubic-unmark.tex
\section{構造対象と $\mathcal O$-marking}\label{sec:weak-definition}
本節では $\kk=\mathbb C$ とする。
非可換 del Pezzo の内在的定義に現れる構造対象の選択を
$\mathcal O$-marking と呼ぶ。
この定義と、後の線束・Weyl 群の定理で用いる追加仮定を分ける。

\begin{definition}[$\mathcal O$-marked del Pezzo と cubic]
\label{def:o-marked-cubic}
$d\in\{1,\ldots,9\}$ とする。
$\mathcal A$ は $\kk$ 線型 noetherian Abel 圏で、$D^b(\mathcal A)$ は
Hom 有限で Serre 関手 $\mathsf S$ を持ち、
$H=\mathsf S^{-1}[2]$ が $H(\mathcal A)=\mathcal A$ を満たすとする。
$\mathcal A_0$ を有限長対象の Serre 部分圏、
$\mathcal A_1$ を $\mathcal A/\mathcal A_0$ で有限長になる対象の逆像とし、
$q:\mathcal A\to\mathcal A/\mathcal A_1$ を商関手とする。
対象 $\OO\in\mathcal A$ が\emph{次数 $d$ の構造対象}であるとは、次を満たすことをいう。
\begin{enumerate}
\item $q\OO$ は非零単純対象であり、
      $\Hom(T,\OO)=0$ がすべての $T\in\mathcal A_1$ に対して成り立つ。
      これを Gabriel--Krull の意味で rank-one torsion-free と呼ぶ。
\item $(\OO,H)$ は Artin--Zhang 豊富である。
      すなわち全対象は十分負の $H$-twist の有限直和の商となり、
      任意の全射 $M\twoheadrightarrow N$ に対して
      $\Hom(\OO,H^nM)\to\Hom(\OO,H^nN)$ は $n\gg0$ で全射となる。
\item 反標準切断環 $R_{\OO}=\bigoplus_{n\ge0}\Hom(\OO,H^n\OO)$ は
      次の Hilbert 級数を持つ。
      \begin{equation}\label{eq:o-del-pezzo-hilbert}
       h_{R_{\OO}}(t)=\frac{1+(d-2)t+t^2}{(1-t)^3},\qquad
       \dim\Hom(\OO,H^n\OO)=1+\frac d2n(n+1)\quad(n\ge0).
      \end{equation}
\end{enumerate}
組 $(\mathcal A,\OO)$ を\emph{$\mathcal O$-marked nc del Pezzo of degree $d$}、
$d=3$ の場合を\emph{$\mathcal O$-marked nc cubic}と呼ぶ。
基礎圏 $\mathcal A$ がこのような対象を少なくとも一つ持つことを、
本稿で採用する無標識 nc del Pezzo の内在的条件とする。
構造対象の例外性、smooth proper dg enhancement、数値 Picard 格子、
AS 表示および ambient 非可換射影空間は、この定義には含めない。
\end{definition}

ここで ``structure sheaf'' はテンソル単位を意味しない。
対象間のテンソル積を仮定せず、同じ $\mathcal A$ の中に複数の構造対象があってよい。
Serre 関手は圏から定まり、$H$ は別の自由な標識データではない。
exact 同値は Gabriel--Krull filtration と Serre 双対性を輸送するので、
上の条件は $(\mathcal A,\OO)$ の exact 同値で不変である。

\begin{lemma}[最小公理から従う消滅と次元]
\label{lem:o-minimal-consequences}
定義\ref{def:o-marked-cubic}の下で、$\mathcal A$ の homological 次元は $2$ である。
また
\begin{equation}\label{eq:o-minimal-vanishing}
 \End(\OO)=\kk,\qquad
 \Hom(\OO,H^{-b}\OO)=0\ (b>0),\qquad
 \Ext^2(\OO,H^n\OO)=0\ (n\ge0).
\end{equation}
$\mathcal A/\mathcal A_1$ の全対象は有限長であり、その単純対象は
$q\OO$ の一つの同型類だけである。
\end{lemma}
\begin{proof}
Serre 双対性から
$\Ext^i(M,N)\simeq\Hom(HN,M[2-i])^*$ となり、$i>2$ では零である。
非零 $M$ に対して $N=H^{-1}M$ と取ると
$\Ext^2(M,N)\simeq\End(M)^*\ne0$ なので次元はちょうど $2$ である。
$\End(\OO)=\kk$ は Hilbert 関数の $n=0$ からも従う。

rank-one torsion-free 対象間の非零射 $f:E\to F$ は単射である。
実際、$qf=0$ なら $\operatorname{im}f\in\mathcal A_1$ が $F$ に埋め込まれて矛盾し、
$qf\ne0$ なら単純対象間の同型なので $\ker f\in\mathcal A_1$ となって核は零である。
$H$ は exact だから、この議論は $H$-twist にも適用できる。
もし $\OO\to H^{-b}\OO$ が非零なら単射であり、
$h(n)=\dim\Hom(\OO,H^n\OO)$ に対して $h(n)\le h(n-b)$ が全整数 $n$ で成り立つ。
各剰余類 $n\bmod b$ でこれは $h(n)$ の有界性を与え、
式\eqref{eq:o-del-pezzo-hilbert}の二次増大と矛盾する。
第二 Ext の消滅は Serre 双対性とこの負の Hom 消滅による。

$j\ge i$ なら $\Hom(H^i\OO,H^j\OO)\ne0$ であり、
その非零射は商圏の単純対象間の同型を与える。
従って全 $qH^i\OO$ は $q\OO$ に同型である。
豊富性による有限生成全射を $q$ で写せば、任意の $qM$ は
$q\OO$ の有限直和の商となる。よって結論を得る。
\end{proof}

\begin{lemma}[中心と自然自己同型]\label{lem:dimension-centre}
定義\ref{def:o-marked-cubic}の下で
\[
 \operatorname{Nat}(\mathrm{id}_{\mathcal A},\mathrm{id}_{\mathcal A})=\kk.
\]
恒等関手の可逆自然変換は $\kk^\times$ であり、$\OO$ 上の恒等を要求すれば自明である。
同じことは、構造対象上の同型を指定した exact 同値の自然自己同型にも成り立つ。
\end{lemma}
\begin{proof}
自然変換 $\eta$ は各 $H^n\OO$ 上でスカラー $c_n$ として作用する。
$\Hom(H^i\OO,H^j\OO)\ne0$ ($j\ge i$) に対する自然性から $c_i=c_j$ である。
任意の対象は有限個の負 twist の直和の商なので、$\eta$ は全対象上で同じスカラーである。
exact 同値の場合はその準逆を合成し、恒等関手の場合に帰着する。
\end{proof}

\begin{proposition}[構造対象による反標準表示]
\label{prop:o-anticanonical-reconstruction}
$\mathcal O$-marked cubic $(\mathcal A,\OO)$ に対して、
Artin--Zhang 再構成は
\[
 (\mathcal A,\OO,H)\simeq(\qgr R_{\OO},\pi R_{\OO},(1))
\]
という exact 同値を与える。
従ってここで採用する構造対象は、反標準表示において実際に
$\pi R_{\OO}$ という構造層として現れる。
\end{proposition}
\begin{proof}
$\OO$ は noetherian、$\End(\OO)=\kk$ は右 noetherian、
全 $\Hom(\OO,M)$ は有限次元であり、$(\OO,H)$ は豊富である。
これらは Artin--Zhang\cite[Theorem 4.5]{AZ}の再構成の仮定である。
\end{proof}

\subsection{線束・Weyl 群の比較で用いる強化条件}
\begin{definition}[例外的構造対象と比較用の弱公理]\label{def:weak-cubic}
定義\ref{def:o-marked-cubic}の $d=3$ の組 $(\mathcal A,\OO)$ が、さらに
$D^b(\mathcal A)$ の smooth proper dg enhancement と
\begin{equation}\label{eq:exceptional}
 \Ext^i(\OO,\OO)=0\qquad(i>0)
\end{equation}
を持つとき、本稿の\emph{比較用の弱公理}を満たすと呼ぶ。
後の数値分類でいう許容線束は、この例外的な構造対象である。
この範囲では rank 条件と Hilbert 関数を
\begin{gather}
 q\OO\text{ が非零単純},\qquad
 \Hom(T,\OO)=0\quad(T\in\mathcal A_1),\label{eq:gk-rank}\\
 \dim\Hom(\OO,H^m\OO)=
 \begin{cases}1+\frac32m(m+1)&m\ge0,\\0&m<0\end{cases}
 \label{eq:weak-hilbert}
\end{gather}
と記す。負方向の消滅は補題\ref{lem:o-minimal-consequences}の帰結である。
以下で「弱標識」と記す組は、この比較用の弱公理を満たす
$\mathcal O$-marking を意味し、最小定義全体を指さない。
\end{definition}

最小の $\mathcal O$-marking とこの強化を同一視しない。
前者の公理から $\Ext^1(\OO,\OO)=0$ が従うかは、本稿では未解決である。
後の命題\ref{prop:o-marking-defect-comparison}は、既知の三次曲面型 heart の中で
この差を数値的欠損と自己拡大によって正確に記述する。
$Q$-marking の先頭対象は例外対象なので、その忘却像は例外的な
$\mathcal O$-marked 部分に入る。

\begin{theorem}[固定圏における構造対象の選択と同値類]
\label{thm:forgetful-fibres}
固定した $\mathcal A$ の次数 $3$ の構造対象の同型類集合を
$\operatorname{Str}_3(\mathcal A)$ とする。
基礎圏との同定を保持した $\mathcal O$-marking の選択は
$\operatorname{Str}_3(\mathcal A)$、同定を忘れた組の exact 同値類は
\[
 \operatorname{Aut}_{\rm ex}(\mathcal A)\backslash
 \operatorname{Str}_3(\mathcal A)
\]
である。例外的部分に限れば $\operatorname{Str}_3$ を
$\operatorname{Line}_{\rm adm}$ に置き換える。
固定した構造対象上の $Q$-marking の選択も同様に対象の同型類で考えるが、
その moduli fiber には構造対象を保つ自己同値の作用が入る。
その群oidと軌道集合の正確な区別は第\ref{sec:marked-moduli}節で与える。
\end{theorem}
\begin{proof}
基礎圏を同定すると構造対象をその圏内に輸送できる。
二つの同定の差は exact 自己同値なので、同定を忘れるとその軌道で同一視する。
例外的部分は許容線束の定義そのものである。
\end{proof}

\section{AOU 関係式モジュライから無標識 Abel 圏へ}\label{sec:aou-unmarked}
本節は関係式の GIT 商と圏の同値類のモジュライを区別し、
両者の間に存在する写像の定義域とファイバーを述べる。
\subsection{三つのモジュライ問題と射の定義域}
AOU\cite[\S5.1, (5.1)--(5.3)]{AOU} の九頂点・十八矢印の有向グラフを $Q$ とする。
関係式のパラメータ空間および矢の基底変換群は
\[
 R_Q:=\prod_{i,j=0}^2\PP^2\simeq(\PP^2)^9,
 \qquad G_Q:=(\mathbb G_m)^{Q_1}\simeq(\mathbb G_m)^{18}.
\]
AOU の関係式 stack と、線形化 $\chi$ に関する射影 GIT 商はそれぞれ
\begin{equation}\label{eq:aou-stack-vs-git}
 \mathfrak R_Q=[R_Q/G_Q],\qquad
 M_Q^{\chi}=R_Q^{\chi\text{-ss}}\mathbin{/\!/}G_Q.
\end{equation}
これらは異なる対象である。後者はトーリック射影スキームであり、
その稠密トーラス $M_Q^\circ\simeq(\mathbb G_m)^8$ を持つ
\cite[Introduction, \S5.4]{AOU}。
特に GIT 商の点が必ず AS 正則 $\mathbb Z$-代数を与えるとは
AOU の結果からは従わない。

\begin{definition}[圏的に許容される AOU 表示]\label{def:aou-categorical-domain}
$T$ 上の\emph{圏的に許容される AOU 表示}とは、AOU の関係式の
$G_Q$-torsor と同変な関係式の族に加え、次を満たすデータである：
(a) 付随する $\{0,1,2\}\times\mathbb Z$-代数から、
$T$ 上で底変更と両立する noetherian Abel 圏 $\mathcal A_T$ が構成される；
(b) 標準九対象 thread は $\mathcal A_T$ の相対完全強例外 helix を誘導し、
その先頭は定義\ref{def:weak-cubic}の相対許容構造対象である；
(c) 相対 Serre 関手、heart、perfect Hom の底変更は
定理\ref{thm:formal-structure}の直前に指定した条件を満たす。
これらの証拠と同型を含めた fibered groupoid を
$\mathfrak R_Q^{\rm cat}$ と書く。
\end{definition}
この定義は $\mathfrak R_Q^{\rm cat}$ が $\mathfrak R_Q$ の代数的な
開部分 stack であることを主張しない。各条件の開性、族としての
構成可能性と降下を検証することが今後の主要課題である。

\begin{proposition}[圏的 AOU 忘却の構成]\label{prop:aou-categorical-map}
上記の相対データが指定されている範囲で、
\begin{equation}\label{eq:aou-categorical-map}
 \mathfrak R_Q^{\rm cat}\longrightarrow\Cw\longrightarrow\Cu,
 \qquad (I,(E_i))\longmapsto(\mathcal A,E_0)
 \longmapsto\mathcal A
\end{equation}
という fibered groupoid の射が定義される。
固定した幾何学的対象 $\mathcal A$ の上で、標識の選択の同型類を
無標識で同一視する関係は、AOU 表示の間に
\emph{$t$-exact な Abel 同値}が存在すること、すなわち
その誘導する圏が exact に同値であることである。
導来同値だけではこの関係を判定できない。
定理\ref{thm:forgetful-fibres}の homotopy quotient の記述は
この射に適用される。
\end{proposition}
\begin{proof}
$G_Q$ は矢の基底を変換するので、relation algebra の頂点を保つ
同型を誘導し、指定された helix と Abel 圏を輸送する。
従って関係式の torsor 上での構成は同型に関して擬自然であり、
定義の相対降下を用いて式\eqref{eq:aou-categorical-map}を得る。
基礎圏の同定を変えるときの同値関係は定理\ref{thm:forgetful-fibres}
そのものである。導来同値が任意の heart を保存するわけではないので、
導来圏のみの比較で Abel 同値を結論することはできない。
\end{proof}

\begin{proposition}[GIT 商への降下に必要な条件]\label{prop:aou-git-descent-test}
関係式の半安定 locus $R_Q^{\chi\text{-ss}}$ の部分集合 $V$ 上で
圏の同型類を割り当てる操作が $M_Q^\chi$ からの写像として
因子化するなら、$V$ の二点が同じ GIT 商点に写るとき、
それらの対応する無標識 Abel 圏は exact に同値でなければならない。
特に S-同値な半安定点に異なる圏が現れれば、その領域には
所望の GIT 商からの忘却写像は存在しない。
\end{proposition}
\begin{proof}
同じ商点の像は一つなので、圏の同型類は一致しなければならない。
これは必要条件のみであり、十分な相対降下を述べない。
\end{proof}

\subsection{可換 locus での完全な比較}
\begin{theorem}[AOU 可換 locus の無標識像]\label{thm:aou-commutative-quotient}
$\mathcal M_{\rm cubic}^{\rm m}$ を滑らかな三次曲面の古典的標識付き
モジュライとする。AOU\cite[Theorem 1.2]{AOU} の局所閉埋込み
\[
 \iota:\mathcal M_{\rm cubic}^{\rm m}\hookrightarrow M_Q^\circ
\]
に沿って式\eqref{eq:aou-categorical-map}を考える。
二つの幾何学的点が同じ無標識 Abel 圏を定めるための必要十分条件は、
対応する滑らかな三次曲面が同型であることである。
従ってこの locus 上の幾何学的同型類の商は、
古典的な $W(E_6)$ による標識変更の商である。
特に $\Aut(X)=1$ の $X$ の上には
\[
 |W(E_6)|=51840
\]
個の古典的標識がある。一般の $X$ では、標識の同型類の個数は
$51840/|\Aut(X)|$ である。
ここでの $W(E_6)$ の作用は\emph{可換 locus 上}の古典的作用であり、
$M_Q^\chi$ 全体に正則作用が延長するとは主張しない。
\end{theorem}
\begin{proof}
AOU の埋込みが与える関係式は、標準九線束の
End 代数を復元し、その Abel 圏は $\operatorname{coh}X$ である。
滑らかな射影多様体の coherent sheaf 圏の exact 同値から
多様体そのものを復元できるので、無標識像が一致することと
$X$ の同型性は同値である\cite{Gabriel}。
古典的標識は六本の順序付き斜交直線の選択と同値であり、
固定 $X$ への標識全体は $W(E_6)$-torsor である
\cite[\S2.3]{AOU}。$\Aut(X)$ は標識全体に自由に作用する：
標識を固定する自己同型は六本の例外線を個別に保つので、
blow-down 後の一般位置の六点を固定する $\PGL_3$ の元であり、恒等である。
したがって標識の同型類数は上記の通りである。
\end{proof}
\begin{remark}[stack と集合の相違]
定理\ref{thm:aou-commutative-quotient}の数は幾何学的同型類集合の
ファイバー数であり、慣性を忘れている。古典的な marked stack を使えば、
可換 locus の無標識 stack は標識付き stack の $W(E_6)$ 商として記述される。
これを非可換 locus に無条件に拡張することはできない。
\end{remark}

\subsection{非可換方向に現れる無限の忘却ファイバー}
\begin{theorem}[一次 AOU helix 選択の無限性：相対仮定付き]\label{thm:aou-infinite-relative-fibre}
$\Aut(X)=1$ の滑らかな三次曲面 $X$ と、
定理\ref{thm:very-general-picard-kernel}の非常に一般の条件を満たす
Poisson 成分を持った許容相対一次 Abel 変形 $\mathcal A_\varepsilon$
を取る。各線束 $L\in\Pic(X)$ と各古典的標準 AOU helix
$E_\bullet$ に対し、$E_\bullet\otimes L$ の許容相対 helix
への持上げが存在するという定理\ref{thm:formal-weyl}の
相対持上げ条件の下では、
\[
 \pi_0\operatorname{Fib}_{\mathcal A_\varepsilon}
 (\mathfrak M_{\rm AOU}^{\rm adm}\longrightarrow\Cu)
 \longrightarrow \Pic(X)/\mathbb ZK_X\simeq\mathbb Z^6
\]
は全射である。特にこの\emph{圏的に許容される helix}
の忘却のファイバーは、同型類の段階でも無限である。
この結論から AOU の射影 GIT 商の幾何学的ファイバーが無限であるとは
直ちには従わない。
\end{theorem}
\begin{proof}
古典的 $E_\bullet\otimes L$ は完全強例外 helix で、先頭は $L$ である。
仮定した相対持上げを取ると、先頭は
定理\ref{thm:formal-structure}の $L_\varepsilon$ に一致する。
一方、系\ref{cor:first-order-pointings}により、
固定 $\mathcal A_\varepsilon$ 上のこれら先頭対象の
exact 自己同値軌道は $\Pic(X)/\mathbb ZK_X$ で分類される。
従って異なるこの商の元に属する二つの helix が
無標識化のファイバー内で同型となることはない。
全ての元に helix の持上げを取れるので全射である。
なお、関係式の GIT 半安定性およびその商への降下は仮定していない。
\end{proof}

\begin{proposition}[無標識慣性の非準コンパクト性に対する障害]\label{prop:nonqc-inertia}
$\operatorname{coh}X$ の exact 自己同値を自然同型で割った
群は $\Pic(X)\rtimes\Aut(X)$ であり、その離散的部分は
$\Pic(X)\simeq\mathbb Z^7$ を含む。
仮に $\Cu$ を正確に表す locally finite type の代数的
$1$-stack があり、$\operatorname{coh}X$ における自己同値の
変形関手が Hochschild の $HH^1(X)$ によって忠実に制御されるなら、
その stack の対角射は準コンパクトではない。
\end{proposition}
\begin{proof}
定理\ref{thm:commutative-structure}によって自己同値の同型類群は
$\Pic(X)\rtimes\Aut(X)$ であり、$\Pic(X)\simeq\mathbb Z^7$。
命題\ref{prop:hh}から $HH^1(X)=0$ である。
忠実な変形制御の仮定により、スカラーの $2$-自己同型を除いた
$1$-stack の幾何学的安定化群は零接空間を持つ。
もし対角射が準コンパクトなら、局所有限型性から
その安定化群は有限型で零次元、従って幾何学的点は有限個である。
しかし $\mathbb Z^7$ の異なる線束捩れは無限個の異なる点を与えるので矛盾する。
\end{proof}
\begin{remark}
この命題は無標識 stack の代数性そのものを否定しない。
また有限型の粗モジュライ空間の存在を否定するものでもない。
対象が持つ exact 自己同値を完全に記録する stack で、
準コンパクトな対角射まで期待してはならないことを示す
\emph{必要条件}である。
\end{remark}

\subsection{大域的認識問題の優先順位}
以上から残る問題は、(i) AOU のどの relation/potential 族で
AS 正則性と相対 heart が保たれるか、(ii) その族の圏への写像の
本質像と exact 同値類、(iii) 自己同値による作用の相対的構成、
(iv) GIT の S-同値との両立、である。
$W(E_6)$ は可換 locus と通常の標識の形式持上げを支配するが、
非可換方向における\emph{無標識化全体の商群}とは同一視できない。
% !TeX root = ../nc-cubic-unmark.tex
\section{数値的線束類の実現と弱標識}\label{sec:line-classification}
固定した Abel 圏の弱標識忘却を理解するためには、まずその圏の中の
線束を分類し、次に exact 自己同値による軌道を求める必要がある。
本節では rank-one torsion-free exceptional 対象を\emph{線束対象}と呼ぶ。
その同型類集合を $\operatorname{Line}(\mathcal A)$、
豊富性と式\eqref{eq:weak-hilbert}も満たす部分を
$\operatorname{Line}_{\rm adm}(\mathcal A)$ と書く。
対象同士のテンソル積も、この集合上の群法則も仮定しない。

以下の主張は固定した反標準 heart の公理から述べる。
数値的線束類の存在は、一意性とは別に証明する必要がある。
内在的線束判定の新規性については網羅的な優先権確認を行っていない。

\subsection{例外性は自己 Euler 形式で判定できる}
以下、定義\ref{def:weak-cubic}を満たす $(\mathcal A,\OO_0)$ を固定する。
$H$ は exact なので $\mathcal A_1$ と rank-one torsion-free 性を保つ。

\begin{lemma}[階数一の射]\label{lem:rank-one-monomorphism}
rank-one torsion-free 対象 $E,F$ の間の非零射は単射である。
また $\End(E)=\mathbb C$ である。
\end{lemma}
\begin{proof}
非零射 $f:E\to F$ に対して $qf=0$ なら、$\operatorname{im}f$ は
$\mathcal A_1$ に属して $F$ に埋め込まれ、torsion-free 性に矛盾する。
従って $qf$ は非零単純対象間の同型である。
$\ker f\in\mathcal A_1$ が $E$ に埋め込まれるので $\ker f=0$ となる。
同じ理由で $\End(E)\to\End(qE)$ は単射である。
右辺は斜体、左辺は Hom 有限性により有限次元の複素代数であるから、
左辺は零因子を持たない有限次元代数である。
任意の元の最小多項式は $\mathbb C$ 上で一次式に分解し、
零因子がないことからその元はスカラーとなる。
\end{proof}

\begin{theorem}[内在的な線束判定]\label{thm:intrinsic-line-criterion}
任意の rank-one torsion-free 対象 $E$ と整数 $b>0$ に対して
\[
 \Hom(E,H^{-b}E)=0,\qquad \Ext^2(E,E)=0,
 \qquad \dim\Ext^1(E,E)=1-\chi(E,E).
\]
特に
\begin{equation}\label{eq:intrinsic-line-criterion}
 E\text{ が線束対象}\quad\Longleftrightarrow\quad\chi(E,E)=1.
\end{equation}
この主張には $E$ 自身の Artin--Zhang 豊富性を仮定しない。
\end{theorem}
\begin{proof}
$\OO_0$ の豊富性から有限直和による全射
$\bigoplus_i H^{-a_i}\OO_0\twoheadrightarrow E$ がある。
非零な成分を選ぶと、前補題により $H^{-a}\OO_0\hookrightarrow E$ を得る。
$h_E(n)=\dim\Hom(\OO_0,H^nE)$ と置けば、$n\ge a$ に対し
\[
 h_E(n)\ge 1+\tfrac32(n-a)(n-a+1).
\]
もし $E\to H^{-b}E$ が非零なら、それも単射であり、$H$ の exact 性から
$h_E(n)\le h_E(n-b)$ がすべての $n$ で成り立つ。
各剰余類 $n\bmod b$ で $h_E(n)$ は上に有界となり、上の二次増大と矛盾する。
Serre 双対性から
$\Ext^2(E,E)=\Hom(E,H^{-1}E)^*=0$ である。
補題\ref{lem:dimension-centre}によって高次 Ext は消え、
前補題の $\dim\End(E)=1$ と Euler 形式の定義から結論を得る。
\end{proof}

\begin{remark}[rank-one だけでは足りない]
可換な滑らかな三次曲面では $E=L\otimes I_Z$ に対し
$\chi(E,E)=1-2\operatorname{length}Z$、
$\dim\Ext^1(E,E)=2\operatorname{length}Z$ となる。
特に点のイデアル層 $I_p$ は rank-one torsion-free だが例外的でない。
一方、定義\ref{def:weak-cubic}の他の条件を保持するなら、
例外性の条件は $\chi(\OO,\OO)=1$ に置き換えてよい。
上の証明が使う $\OO_0$ の条件は rank-one、豊富性と切断数の増大であり、
この置換の確認に例外性を先取りする必要はない。
\end{remark}

\subsection{数値的欠損と同一数値類での一意性}
次の三次曲面型の数値的構造は\emph{追加仮定}である。
$K_0^{\rm num}(\mathcal A)$ に三次曲面型の表示
$\operatorname{ch}(E)=(r,D,s)$、$r\in\mathbb Z$、$D\in\Lambda=I_{1,6}$、
$K^2=3$ を指定し、$s-D^2/2\in\mathbb Z$、
\begin{align}
 \chi((r,D,s),(r',D',s'))
 &=rr'+rs'+r's-D\cdot D'
   -\tfrac12 K\cdot(rD'-r'D),\label{eq:line-rr}\\
 \operatorname{ch}(HE)
 &=(r,D-rK,s-D\cdot K+\tfrac12rK^2)\label{eq:line-H}
\end{align}
をすべての対象について仮定し、$r(\OO_0)=1$ と正規化する。
数値的階数と内在的 rank 条件の対応、および torsion 対象の次数正値性は、
以下でこの仮定から導く。
数値的線束類とは $v_D=(1,D,D^2/2)$ のことである。

\begin{theorem}[反標準 Serre 消滅と内在的階数]\label{thm:ac-serre-positivity}
定義\ref{def:weak-cubic}と式\eqref{eq:line-rr}--\eqref{eq:line-H}、
$r(\OO_0)=1$ を仮定する。$d(M)=(-K)\cdot c_1(M)$、$d_0=d(\OO_0)$ と置く。
任意の $M\in\mathcal A$ に対して
\begin{equation}\label{eq:ac-serre-vanishing}
 \Ext^i(\OO_0,H^nM)=0\qquad(i>0,\ n\gg0)
\end{equation}
であり、十分大きな $n$ で
\begin{equation}\label{eq:ac-hilbert-polynomial}
 \dim\Hom(\OO_0,H^nM)
 =\tfrac32r(M)n(n+1)+(d(M)-r(M)d_0)n+\chi(\OO_0,M).
\end{equation}
さらに次が成り立つ。
\begin{enumerate}
\item 全対象で $r(M)\ge0$。$r(M)=0$ なら $d(M)\ge0$。
\item $r(M)=d(M)=0$ なら $M\in\mathcal A_0$。$M\ne0$ なら
      $\chi(\OO_0,M)>0$ であり、$\operatorname{length}M\le\chi(\OO_0,M)$。
\item $\mathcal A_1=\{M:r(M)=0\}$。rank-one torsion-free 対象 $E$ には $r(E)=1$。
\end{enumerate}
特に有効次数正値性
\begin{equation}\label{eq:effective-degree-positivity}
 d(T)\ge0\qquad(T\in\mathcal A_1) \tag{EP}
\end{equation}
は独立の仮定を要しない。
\end{theorem}
\begin{proof}
$D_0=c_1(\OO_0)$ とする。$\OO_0$ の例外性と Riemann--Roch より
$\operatorname{ch}(\OO_0)=v_{D_0}$ であり、
$\chi(\OO_0,H^m\OO_0)=1+3m(m+1)/2$ である。
$m\ge0$ では Serre 双対性と所定の負の Hom 消滅から
$\Ext^2(\OO_0,H^m\OO_0)=0$ である。
式\eqref{eq:weak-hilbert}との差を取れば
$\Ext^1(\OO_0,H^m\OO_0)=0$ も得る。

豊富性により $P=\bigoplus_jH^{-a_j}\OO_0\twoheadrightarrow M$ を選ぶ。
この全射は
\[
 \Hom(M,H^{-n-1}\OO_0)
 \hookrightarrow\bigoplus_j\Hom(\OO_0,H^{a_j-n-1}\OO_0)
\]
を与える。右辺は $n\gg0$ で零なので、Serre 双対性から
$\Ext^2(\OO_0,H^nM)=0$ となる。
$N=\ker(P\to M)$ にもこの第二 Ext 消滅を適用すると、
長完全列の
\[
 \Ext^1(\OO_0,H^nP)\longrightarrow\Ext^1(\OO_0,H^nM)
 \longrightarrow\Ext^2(\OO_0,H^nN)
\]
の両端が $n\gg0$ で零となる。
高次 Ext は補題\ref{lem:dimension-centre}で消える。
これで式\eqref{eq:ac-serre-vanishing}を得て、Riemann--Roch を展開すれば
式\eqref{eq:ac-hilbert-polynomial}となる。

非零対象 $M$ ではこの多項式は恒等的に零でない。
実際、AZ 豊富性は任意に負の tail 内の有限直和による生成を与える
\cite[Lemma 2.1(ii)]{PolishchukProj}ので、
$\Hom(\OO_0,H^nM)$ は任意に大きい $n$ のうち少なくとも一つで非零である。
非負の値を取る非零多項式の最高次係数を順に見ると、(1)と、
$r(M)=d(M)=0$, $M\ne0$ のときの $\chi(\OO_0,M)\in\mathbb Z_{>0}$ を得る。

$r(M)=d(M)=0$ のとき、加法性と (1) により全部分商の $r,d$ も零である。
従って任意の有限な厳密部分対象列の各非零商に対して
$\chi(\OO_0,-)$ は正整数であり、列の長さは $\chi(\OO_0,M)$ 以下である。
これが (2) を示す。

次に $r(M)=0$ とし、$Q_0:\mathcal A\to\mathcal A/\mathcal A_0$ と置く。
$Q_0M$ の有限な厳密部分対象列は、Serre 商の部分対象の持上げにより
$M$ 内の有限な部分対象列で表せる。必要なら代表の累積和を取り、
始点を $0$、終点を $M$ とする。
各厳密段階の商 $N$ は $r(N)=0$ かつ $Q_0N\ne0$ だから、(2) より $d(N)>0$ である。
$d(N)$ は整数なので、この列の長さは $d(M)$ 以下である。
従って $Q_0M$ は有限長、すなわち $M\in\mathcal A_1$ となる。
この議論は $d$ が $\mathcal A_0$ 上で零になることを用いていない。

逆に $T\in\mathcal A_1$ に対し、全射
$\bigoplus_jH^{-a_j}\OO_0\twoheadrightarrow T$ の各成分の像を $I_j$ とする。
$\mathcal A_1$ は Serre 部分圏なので $I_j\in\mathcal A_1$ であり、
(1)と $r(\OO_0)=1$ から $r(I_j)\in\{0,1\}$ である。
$r(I_j)=1$ なら、その成分の核の階数は零、従って既に示した向きにより
核も $\mathcal A_1$ に入る。すると $H^{-a_j}\OO_0\in\mathcal A_1$ となり、
$q\OO_0\ne0$ と矛盾する。
従って全 $I_j$ の階数は零であり、その有限和 $T$ も階数零である。
これで $\mathcal A_1=\{r=0\}$ を得る。

rank-one torsion-free 対象 $E$ には豊富性と
補題\ref{lem:rank-one-monomorphism}から $H^{-a}\OO_0\hookrightarrow E$ が存在する。
商圏で非零単純対象間の同型になるので、その余核は $\mathcal A_1$ に入る。
従って $r(E)=1$。最後に (EP) は (1)と (3) の帰結である。
\end{proof}

\begin{remark}[入力と結論の範囲]\label{rem:serre-inputs}
上の消滅は一般の AZ 豊富性だけから導いたものではない。
指定した Serre 関手、三次曲面型 Riemann--Roch、構造対象の Hilbert 関数を使っている。
生成には異なる指数の有限直和だけを使い、各 $n\gg0$ について
単一の $H^{-n}\OO_0$ の直和で生成できるとは主張しない。
また (2) の逆、すなわち $M\in\mathcal A_0\Rightarrow d(M)=0$ もここでは未証明である。
\end{remark}

\begin{corollary}[零次元欠損の数値的類似]\label{cor:line-numerical-defect}
rank-one torsion-free 対象 $E$ の表示を $(1,D,s)$ とし、
$n(E)=D^2/2-s$ と置く。このとき
\[
 n(E)\in\mathbb Z_{\ge0},\qquad
 \chi(E,E)=1-2n(E),\qquad
 \dim\Ext^1(E,E)=2n(E).
\]
従って数値的線束類を持つ rank-one torsion-free 対象は自動的に例外的である。
\end{corollary}
\begin{proof}
式\eqref{eq:line-rr}に $E$ を二度代入し、
定理\ref{thm:intrinsic-line-criterion}を適用する。
\end{proof}

一意性には自己 Ext だけではなく二対象間の Ext が必要である。
以下の Serre 双対性と次数による消滅論法は、異なる設定である
terminal del Pezzo order 上の先行例\cite[Lemma 3.9]{HR}とも共通する。
そのため定理\ref{thm:ac-serre-positivity}で導いた (EP) を用いる。
有限代数の Euler 行列だけの一致は、この定理の反標準 heart の仮定を代用しない。

\begin{theorem}[数値的線束類における一意性]\label{thm:numerical-line-uniqueness}
定理\ref{thm:ac-serre-positivity}の数値仮定を置く。
rank-one torsion-free 対象 $E,F$ が $d(E)<d(F)+3$ を満たせば
$\Ext^2(E,F)=0$ である。特に $[E]=[F]=v_D$ なら $E\cong F$ である。
従って $\operatorname{Line}(\mathcal A)$ から $\Lambda$ への第一 Chern 類写像は単射である。
\end{theorem}
\begin{proof}
Serre 双対性により $\Ext^2(E,F)=\Hom(F,H^{-1}E)^*$ である。
右辺に非零射があれば、補題\ref{lem:rank-one-monomorphism}により単射であり、
余核 $T$ は $\mathcal A_1$ に属す。
その次数は式\eqref{eq:line-H}から
\[
 d(T)=d(H^{-1}E)-d(F)=d(E)-3-d(F)<0
\]
となり (EP) に矛盾する。
$E,F$ が同じ数値的線束類なら両向きの第二 Ext は零であり、
$\chi(E,F)=\chi(F,E)=1$ なので両向きに非零射が存在する。
いずれも単射であり、合成は非零である。
$\End(E)=\mathbb C$ だから合成は可逆なスカラーとなる。
一方の単射は左逆を持ち、他方も単射であることから互いに同型である。
最後に任意の線束対象の数値類は前系により $v_D$ なので、単射性が従う。
\end{proof}

\begin{remark}[存在・許容性・軌道を分ける]\label{rem:line-realization-set}
実現される因子類の集合を $\Omega_{\mathcal A}\subseteq\Lambda$ とすると、
上の仮定の下で $\operatorname{Line}(\mathcal A)\cong\Omega_{\mathcal A}$ である。
$H$ は線束対象を保つので $\Omega_{\mathcal A}-K=\Omega_{\mathcal A}$。
しかし $\Omega_{\mathcal A}=\Lambda$ は存在定理であり、一意性からは従わない。
また $\operatorname{Line}_{\rm adm}(\mathcal A)=\operatorname{Line}(\mathcal A)$ も
一般にはまだ示していない。
実際、線束対象 $E$ に対して $m\ge0$ なら
\[
 \dim\Hom(E,H^mE)
 =1+\tfrac32m(m+1)+\dim\Ext^1(E,H^mE),
\]
ここで第二 Ext の消滅は定理\ref{thm:intrinsic-line-criterion}から従う。
所定の切断数のためには交差 Ext の消滅、豊富性のためには生成条件が別に必要である。
さらに $H$ の軌道商と全 exact 自己同値の軌道商を同一視しない。
\end{remark}

\begin{proposition}[許容性の生成・消滅判定]\label{prop:line-admissibility-test}
定理\ref{thm:ac-serre-positivity}の数値仮定の下で、線束対象 $L$ が許容されることは、
次の二条件と同値である。
\begin{enumerate}
\item 有限和による全射 $\bigoplus_j H^{-a_j}L\twoheadrightarrow\OO_0$ が存在する。
\item $\Ext^{i}(L,H^mL)=0$ が $i>0$, $m\ge0$ に対して成立する。
\end{enumerate}
\end{proposition}
\begin{proof}
$L$ が許容なら第一条件は AZ 生成である。
第二 Ext は定理\ref{thm:numerical-line-uniqueness}で消え、
規定の切断数と Riemann--Roch の比較から第一 Ext も消える。
逆に二条件を仮定する。
第一条件を反標準 twist し、$\OO_0$ の任意 tail による生成と合成すると、
$L$ の任意 tail の有限和も全対象を生成する。
任意の $M$ に対して $P=\bigoplus_jH^{-b_j}L\twoheadrightarrow M$ を取り、
その核を $N$ とする。
負の自己 Hom 消滅と Serre 双対性から
$\Ext^2(L,H^nM)=\Ext^2(L,H^nN)=0$ が $n\gg0$ で成立する。
第二条件は $\Ext^1(L,H^nP)=0$ を与えるから、長完全列により
$\Ext^1(L,H^nM)=0$ も従う。
任意の全射の核にこの消滅を適用すれば AZ の Hom 全射条件を得る。
切断数は第二条件と Riemann--Roch、負方向は
定理\ref{thm:intrinsic-line-criterion}の負の自己 Hom 消滅から得られる。
従って $L$ は許容構造対象である。
\end{proof}

\subsection{零次数単純対象と許容性の障害}
ここでも定理\ref{thm:ac-serre-positivity}の数値仮定を置く。
次の命題は、線束対象であることから AZ 豊富性へ進む際に
調べるべき障害を与える。障害となる単純対象の実在は仮定であり、
本稿はその具体例を構成したとは主張しない。

\begin{lemma}[零次数単純対象の数値類]\label{lem:zero-degree-simple-classification}
$d(T)=0$ なる単純対象 $T$ に対して $C=c_1(T)$ と置く。
このとき $r(T)=0$ であり、$C=0$ または
\[
 C^2=-2,\qquad HT\cong T,\qquad \Ext^1(T,T)=0
\]
のいずれかである。後者では
$c=\dim\Hom(\OO_0,T)=\chi(\OO_0,T)>0$ となる。
\end{lemma}
\begin{proof}
$T\in\mathcal A_0\subset\mathcal A_1$ だから $r(T)=0$ である。
Hom 有限性と単純性により $\End(T)=\mathbb C$。
$H^{-1}T$ も単純なので
\[
 \epsilon:=\dim\Ext^2(T,T)=\dim\Hom(T,H^{-1}T)\in\{0,1\},
 \qquad \dim\Ext^1(T,T)=1+\epsilon+C^2.
\]
$C\in K^\perp$、$K^2=3$ と $\Lambda$ の符号 $(1,6)$ から
$K^\perp$ は負定値である。
また $D_0=c_1(\OO_0)$ とすると
$\chi(\OO_0,T)=s(T)-D_0\cdot C$ は整数だから $s(T)$ は整数であり、
$s(T)-C^2/2\in\mathbb Z$ から $C^2$ は偶数である。
従って $C\ne0$ なら $C^2=-2$、$\epsilon=1$、$\Ext^1(T,T)=0$ となる。
$\epsilon=1$ は $HT\cong T$ を意味する。
この場合、定理\ref{thm:ac-serre-positivity}と $HT\cong T$ から
$\Ext^{>0}(\OO_0,T)=0$ を得る。
同定理の零階数・零次数の場合の正値性により
$c=\chi(\OO_0,T)>0$ である。
\end{proof}

\begin{lemma}[数値的点への Ext 消滅]\label{lem:numerical-point-ext}
$P$ を $\operatorname{ch}(P)=(0,0,s)$ なる単純対象とする。
このとき $s>0$ であり、任意の線束対象 $L$ について
\[
 \Ext^1(P,L)=\Ext^1(L,P)=0,\qquad \dim\Hom(L,P)=s.
\]
\end{lemma}
\begin{proof}
$s=\chi(\OO_0,P)>0$ は定理\ref{thm:ac-serre-positivity}から従う。
非分裂な拡大 $0\to L\to F\to P\to0$ が存在すると仮定する。
$F$ の非零 torsion 部分対象は $L$ と交わらず $P$ に単射するため、
単純性から $P$ と同型になり、拡大を分裂させてしまう。
従って $F$ は rank-one torsion-free である。
しかし $\operatorname{ch}(F)=\operatorname{ch}(L)+(0,0,s)$ なので
$n(F)=-s<0$ となり、系\ref{cor:line-numerical-defect}に矛盾する。
これで $\Ext^1(P,L)=0$ を得る。
$H^{-1}L$ も線束対象であり、Serre 双対性から
$\Ext^1(L,P)\cong\Ext^1(P,H^{-1}L)^*=0$。
さらに $\Ext^2(L,P)=\Hom(P,H^{-1}L)^*=0$ なので、
$\chi(L,P)=s$ が Hom 次元を与える。
\end{proof}

\begin{theorem}[root 単純対象に対する線束の反射]\label{thm:root-line-reflection}
定理\ref{thm:ac-serre-positivity}の仮定を置く。
$T$ を $d(T)=0$, $C=c_1(T)\ne0$ なる単純対象とし、
$\operatorname{ch}(T)=(0,C,s)$ と書く。
任意の線束対象 $L$ に対して $D=c_1(L)$、
$a=\chi(L,T)=s-D\cdot C$ と置くと、
\begin{equation}\label{eq:root-line-hom-ext}
 \dim\Hom(L,T)=\max(a,0),\qquad
 \dim\Ext^1(L,T)=\dim\Ext^1(T,L)=\max(-a,0).
\end{equation}
線束の同型類集合上に involution $\rho_T$ があり、
\begin{equation}\label{eq:root-affine-reflection}
 c_1(\rho_TL)=D-aC=D+(D\cdot C-s)C
\end{equation}
を満たす。$a>0$ では評価全射の核、$a<0$ では普遍拡大、
$a=0$ では $L$ 自身によって実現される：
\[
 \begin{array}{ll}
 a>0:&0\longrightarrow\rho_TL\longrightarrow L
          \longrightarrow T^{\oplus a}\longrightarrow0,\\[2pt]
 a<0:&0\longrightarrow L\longrightarrow\rho_TL
          \longrightarrow T^{\oplus(-a)}\longrightarrow0.
 \end{array}
\]
さらに $R_{\rho_TL}\cong R_L$ である。
\end{theorem}
\begin{proof}
補題\ref{lem:zero-degree-simple-classification}により
$C^2=-2$, $HT\cong T$, $\Ext^1(T,T)=0$ である。
torsion-free 性から $\Hom(T,L)=0$ と
$\Ext^2(L,T)=\Hom(T,H^{-1}L)^*=0$ を得る。
Serre 双対性と $HT\cong T$ より
$\Ext^1(T,L)\cong\Ext^1(L,T)^*$。
$h=\dim\Hom(L,T)$、$b=\dim\Ext^1(T,L)$ と置けば $a=h-b$ である。

$V=\Ext^1(T,L)$ とし、
$\Ext^1(V\otimes T,L)=V^*\otimes V$ の恒等元による拡大
\[
 0\longrightarrow L\longrightarrow F\longrightarrow V\otimes T
 \longrightarrow0
\]
を取る。境界写像 $\Hom(T,V\otimes T)=V\to\Ext^1(T,L)=V$ は同型なので、
$\Hom(T,F)=0$ である。
$F$ の torsion 部分対象は $L$ と交わらず、半単純対象 $V\otimes T$ に埋め込まれる。
非零なら $T$ を部分対象として含み、上の消滅に反する。
従って $F$ は torsion-free であり、$qF\cong qL$ は非零単純なので階数一である。
数値的欠陥を計算すると
\[
 \begin{aligned}
 n(F)
 &=\frac{(D+bC)^2}{2}-\left(\frac{D^2}{2}+bs\right)\\
 &=-b(s-D\cdot C)-b^2=-b(h-b)-b^2=-bh.
 \end{aligned}
\]
系\ref{cor:line-numerical-defect}より $n(F)\ge0$ だから $bh=0$。
$h,b\ge0$ と $a=h-b$ から式\eqref{eq:root-line-hom-ext}が従う。

$a<0$ なら $b=-a$ であり、上の $F$ は $n(F)=0$ なので線束対象である。
$a>0$ なら $h=a$ であり、全 Hom による評価 $L\to T^{\oplus a}$ は全射である。
実際、像が真の部分対象なら、それを殺す非零射 $T^{\oplus a}\to T$ が
$\Hom(L,T)$ の基底間に非自明な関係を与える。
核 $E$ は rank-one torsion-free であり、
\[
 n(E)=a(s-D\cdot C)+\frac{a^2C^2}{2}=a^2-a^2=0
\]
なので線束対象である。これで各場合の $\rho_TL$ を得る。
どちらの場合も式\eqref{eq:root-affine-reflection}が成立し、
$s-(D-aC)\cdot C=-a$ である。
二回反射した Chern 類は $D$ に戻るので、
定理\ref{thm:numerical-line-uniqueness}から $\rho_T^2L\cong L$ を得る。
従って構成時の基底の選択も対象の同型類には影響しない。

環同型は $a>0$ の評価列で証明すればよい。
式\eqref{eq:root-line-hom-ext}より $\Ext^1(L,T)=0$ であり、
評価列から $\Hom(E,T)=0$ となる。
任意の整数 $n$ に対して
\[
 \Hom(T,H^nL)=0,\qquad
 \Ext^1(T,H^nL)\cong\Ext^1(L,H^{-n-1}T)^*=0.
\]
評価列に $\Hom(-,H^nL)$ を適用し、次に
$\Hom(E,H^nT)=0$ を使うと、
\[
 \Hom(L,H^nL)\cong\Hom(E,H^nL)\cong\Hom(E,H^nE)
\]
を得る。これは $L\to H^nL$ を $E$ に制限する写像であり、
$H$ と射の合成に両立するので、$R_L\cong R_E$ を与える。
$a<0$ の場合は反射後の正の側からこの議論を適用し、
$a=0$ の場合は恒等写像を取ればよい。
\end{proof}

\begin{corollary}[許容線束の半空間条件と root 排除]\label{cor:admissible-root-halfspace}
上の $T$ に対し、$L\in\operatorname{Line}_{\rm adm}(\mathcal A)$ なら
\[
 s-c_1(L)\cdot C>0.
\]
特にすべての $D\in\Lambda$ が許容線束の第一 Chern 類として実現されるなら、
零次数単純対象はすべて $c_1=0$ を満たす。
\end{corollary}
\begin{proof}
$a\le0$ なら定理より $\Hom(L,T)=0$ であり、$HT\cong T$ なので
$\Hom(H^{-n}L,T)=0$ がすべての $n$ で成立する。
これは $(L,H)$ の AZ 生成条件に反する。
最後の主張では、$C\ne0$ なる root simple があれば $D=-mC$ を取り、
$s-D\cdot C=s-2m\le0$ となる十分大きい整数 $m$ を選ぶと矛盾する。
\end{proof}

\begin{remark}
式\eqref{eq:root-affine-reflection}は $D\cdot C=s$ を壁とする affine reflection であり、
線束 $L$ 自体を変える。固定弱標識上の有限 Weyl 標識変更とは区別する。
構成は spherical twist の評価・普遍拡大と同じ数値的形式を持つが、
ここでは固定 heart 内の完全列によって証明した。
環同型から $L$ と $\rho_TL$ を交換する exact 自己同値が得られるとは主張しない。
半空間条件は必要条件であり、その正の側にある全線束の豊富性を結論しない。
\end{remark}

\begin{proposition}[同じ切断環を持つ非許容線束]\label{prop:root-simple-obstruction}
零次数 root 型単純対象 $T$ が存在すると仮定し、
$c=\dim\Hom(\OO_0,T)>0$ とする。
評価全射の核
\[
 0\longrightarrow E\longrightarrow\OO_0\longrightarrow T^{\oplus c}
 \longrightarrow0
\]
は線束対象であり、制限写像は次数付き環同型 $R_{\OO_0}\cong R_E$ を与える。
一方 $(E,H)$ は AZ 豊富ではない。
\end{proposition}
\begin{proof}
補題\ref{lem:zero-degree-simple-classification}と
定理\ref{thm:root-line-reflection}を $L=\OO_0$ に適用する。
反射後は $\chi(E,T)=-c$ なので $\Hom(E,T)=0$ であり、
$HT\cong T$ から全ての $n$ に対して $\Hom(H^{-n}E,T)=0$ となる。
従って AZ 生成条件が成立しない。
環同型は同定理の制限写像による証明で与えられる。
この命題は $T$ の存在を仮定するものであり、実例の構成ではない。
\end{proof}


\subsection{AS heart と反標準 heart の比較}
本稿で比較するのは $\mathcal A_{\rm AS}$ と $\mathcal A_{\rm ac}$ だけである。
次の十分条件は、指定した導来同値を exact 同値へ強める際に用いる。
\begin{lemma}[消滅と生成による heart の識別]\label{lem:ample-heart-transfer}
$\mathcal A,\mathcal B$ の有界導来圏の同値 $\Phi$ と、
各 heart 内の列 $P_n,Q_n$ ($n\ge0$) があるとする。
各対象 $M$ に対し十分大きな $n$ で
$\Ext^{p}(P_n,M)=0$ ($p>0$) かつ有限直和 $P_n^{\oplus a}\twoheadrightarrow M$ が存在し、
$Q_n$ についても同じ性質を仮定する。
十分大きなすべての $n$ で $\Phi(P_n)\cong Q_n$ なら $\Phi$ は $t$-exact である。
\end{lemma}
\begin{proof}
$C\in D^b(\mathcal A)$ の有限個の非零 cohomology 対象に対して
一様に大きい $n$ を取る。標準 truncation と Ext 消滅から
\[
 \Hom(P_n,C[i])\cong\Hom(P_n,H^i(C))
\]
となる。生成全射の存在によって右辺がすべての十分大きい $n$ で零であることと
$H^i(C)=0$ は同値である。
従って $C$ が $\mathcal A$ の heart に属すことは、
$i\ne0$ に対する上の Hom 消滅で特徴付けられる。
$\Phi$ とその逆に同じ議論を適用すれば、二つの heart が一致する。
\end{proof}


この補題では、両側の列をそれぞれ指定された構造対象の
十分負の反標準列として取る。両側での消滅・生成と列の同定は、
数値的な $K_0$ の一致や導来同値の存在だけでは代用できない。
AS heart で得た線束を反標準 heart の線束と認識する際も、
このような比較の仮定を明記する。

可換点を特殊ファイバーとする許容相対 Artin 変形では、
各古典的線束の一意な持上げは定理\ref{thm:formal-structure}で与えられる。
全数値類の同時存在は、後の定理\ref{thm:very-general-line-existence}で
明示する代数族の仮定の下では非常に一般の点で証明できる。
その適用には指定された反標準 heart と代数族との比較が必要であり、
任意の非可換点での存在や全域的な開性は主張しない。

% !TeX root = ../nc-cubic-unmark.tex
\section{AOU 型線束標識の有限性と反射への還元}\label{sec:weyl-frame}
\subsection{標準九対象からの格子標識の復元}
以下では AOU の三ブロック型の\emph{線束}標識に対象を限定する。
OU\cite[Definition 1.1(3)]{OU2025}の任意の very strong helix は、
この型に限られないので、以下の有限性をその全体には主張しない。
定義\ref{def:weak-cubic}を満たす $(\mathcal A,\OO_0)$ を固定し、
反標準 heart $\mathcal A$ に式\eqref{eq:line-rr}--\eqref{eq:line-H}を仮定し、
$r(\OO_0)=1$ の正規化を置く。
線束対象の第一 Chern 類写像は定理\ref{thm:numerical-line-uniqueness}により単射である。

$\Lambda_{\rm std}=\mathbb Zh_0\oplus\bigoplus_{i=1}^6\mathbb Ze_i^0$ を
$h_0^2=1$, $(e_i^0)^2=-1$ なる直交基底を持つ格子とし、
$K_0=-3h_0+\sum_i e_i^0$ とする。
標識格子の候補の集合を
\[
 \mathfrak F(\Lambda,K)
 =\operatorname{Isom}((\Lambda_{\rm std},K_0),(\Lambda,K))
\]
とする。これは非空なら前合成による右 $W(E_6)$-torsor である
\cite[\S2.6]{DGK}。

固定線束 $L$ に対し、$Q$ の頂点でラベル付けされた九対象を
\[
 (L,A,B;X_1,X_2,X_3;Y_1,Y_2,Y_3)
\]
と書く。すべて線束対象であり、各ブロックは互いに完全直交し、
この順の Euler 行列が
\begin{equation}\label{eq:frame-euler}
 C_Q=\begin{pmatrix}I_3&J_3&2J_3\\0&I_3&J_3\\0&0&I_3\end{pmatrix},
 \qquad J_3=(1)_{3\times3}
\end{equation}
である標準型を考える。行が射の始域、列が終域である。
本節で $\operatorname{Mark}_Q^{\rm line}(\mathcal A,L)$ と書く集合は、
この Euler 型を持ち、full strong thread から $E_{i+9}=HE_i$ により得る helix が
very strong である九つ組の同型類とする。各対象は指定された反標準 heart に属す。
指定 AS 表示の再構成条件も含めて AOU/OU marking とみなす場合には、
その条件を満たす部分集合に限定する。以下の単射と上界はその部分集合にも成立する。
対象の同型類を数え、矢の基底や自己同値による同一視は含めない。

\begin{lemma}[Euler 行列による標識復元]\label{lem:weyl-frame-reconstruction}
上の九対象の $L$ に対する相対第一 Chern 類を
$(0,a,b;x_1,x_2,x_3;y_1,y_2,y_3)$ と置く。このとき
\begin{equation}\label{eq:frame-reconstruction}
 h=x_1+x_2+x_3-b,\qquad
 e_i=x_i\ (1\le i\le3),\qquad e_{3+j}=h-y_j\ (1\le j\le3)
\end{equation}
は $\Lambda$ の整基底であり、
\[
 h^2=1,\quad e_i^2=-1,\quad h\cdot e_i=e_i\cdot e_j=0\ (i\ne j),
 \qquad K=-3h+\sum_{i=1}^6e_i.
\]
さらに $a=K+h$ と $b=-h+e_1+e_2+e_3$ が成り立つ。
従って九対象は一意な $\varphi\in\mathfrak F(\Lambda,K)$ を定め、
その相対 Chern 類は AOU\cite[(4.1), (9.44)]{AOU}の標準九類の $\varphi$ による像である。
\end{lemma}
\begin{proof}
数値的線束類 $v_D,v_E$ に対する Riemann--Roch から
\[
 \chi(v_D,v_E)=1+\frac{(E-D)^2-K\cdot(E-D)}2
\]
である。従って式\eqref{eq:frame-euler}は
\[
 -K\cdot a=-K\cdot b=0,\quad -K\cdot x_i=1,\quad -K\cdot y_j=2
\]
と、次の交点数を与える：
\[
 \begin{gathered}
 a^2=b^2=-2,\quad a\cdot b=-1,\quad
 x_i^2=-1,\quad x_i\cdot x_j=0\ (i\ne j),\\
 y_i^2=0,\quad y_i\cdot y_j=1\ (i\ne j),\quad
 x_i\cdot y_j=0,\\
 a\cdot x_i=b\cdot x_i=a\cdot y_j=b\cdot y_j=-1.
 \end{gathered}
\]
これを式\eqref{eq:frame-reconstruction}へ代入すると、$h,e_i$ の Gram 行列は
$\operatorname{diag}(1,-1,\ldots,-1)$ になる。その行列式の絶対値が1であり、
$\Lambda$ も階数7の unimodular 格子なので、生成する部分格子の指数は1である。
次数の計算から $K\cdot h=-3$, $K\cdot e_i=-1$ となり、
$K=-3h+\sum_i e_i$ を得る。
$a$ と $K+h$ はこの基底のすべての元との交点数が同じなので等しい。
$b$ の式は $h$ の定義である。
\end{proof}

\begin{theorem}[固定弱標識上の有限性と反射への還元]\label{thm:weyl-frame-injection}
上の仮定の下で、固定 $(\mathcal A,L)$ における標準三ブロック線束標識の集合
$\operatorname{Mark}_Q^{\rm line}(\mathcal A,L)$ から
$\mathfrak F(\Lambda,K)$ への復元写像は単射である。特に
\[
 \#\operatorname{Mark}_Q^{\rm line}(\mathcal A,L)\le51840.
\]
この集合が非空であり、各実現された標識の格子標識を任意の単純反射で
前合成したものも同じ種類の標識として実現されるなら、この写像は全単射である。
この場合、実際の標識集合は $W(E_6)$-torsor となり、
反射操作の Coxeter 関係は対象の同型類上で自動的に成り立つ。
\end{theorem}
\begin{proof}
同じ格子標識を与える二つの標識では、対応する九対象の第一 Chern 類が
$c_1(L)$ を含めて一致する。線束の数値的一意性により、対応する対象は同型である。
helix の残りの対象は指定した $H$ の反復で定まり、矢の基底は数えないので、
二つの標識は同じである。これで単射性を得る。
標識集合が非空なら復元補題により $\mathfrak F(\Lambda,K)$ も非空である。
$\mathfrak F(\Lambda,K)$ は位数51840の $W(E_6)$ の torsor なので上界が従う。
反射に関する存在条件の下では、非空な像はすべての単純反射に安定であり、
単純反射は $W(E_6)$ を生成するから像は全体である。
各反射操作はこの全単射を通じた $W(E_6)$ の作用に等しい。
従って異なる反射語が同じ Weyl 元を与えれば、対応する対象の数値類が一致し、
単射性によって同じ標識を与える。
\end{proof}

\begin{corollary}[一種類の三対象変換への還元]\label{cor:weyl-one-move}
標準型は先頭 $L$ を保つ各直交ブロック内の頂点ラベルの置換に安定であるとする
（上で定義した九つ組の集合では成立し、追加の AS 表示条件を課す場合もその安定性を確認する）。
上の定理の反射閉性を示すには、すべての標識について次の変換を実現すれば十分である：
\begin{equation}\label{eq:weyl-cross-block-move}
 \begin{gathered}
 (0,a,b;x_1,x_2,x_3;y_1,y_2,y_3)
 \longmapsto(0,a,b';x_1,x_2,x_3';y_1',y_2,y_3),\\
 b'=x_1+x_2-y_1,\qquad x_3'=h-y_1,\qquad y_1'=h-x_3.
 \end{gathered}
\end{equation}
ここで必要なのは、指定した相対 Chern 類を持つ三つの線束対象の存在に加え、
置換後の九対象が同じ標準型の marking となることである。
\end{corollary}
\begin{proof}
単純根を
\[
 \alpha_i=e_i-e_{i+1}\ (1\le i\le5),\qquad
 \alpha_6=-h+e_4+e_5+e_6=a-b
\]
と取る。これらの負 Gram 行列は $E_6$ の Cartan 行列であり、
$s_\alpha(v)=v+(v\cdot\alpha)\alpha$ は $K$ を固定する。
現在の標識 $\varphi$ において
$s_{\varphi(\alpha)}\circ\varphi=\varphi\circ s_\alpha$ であるから、
以下の反射は上で定義した右作用と両立する。
$s_1,s_2$ は $X$ ブロックの隣接交換、$s_4,s_5$ は $Y$ ブロックの隣接交換、
$s_6$ は最初のブロック内の $A,B$ の交換である。
いずれも $L$ を保ち、仮定により標識上で実現される。
残る $s_3$ は $e_3,e_4$ を交換し、それ以外の $h,e_i$ を固定する。
標準九類へ代入すれば式\eqref{eq:weyl-cross-block-move}となる。
\end{proof}

ブロック内の置換が very strongness を保つことは、周期を越えた射を含めて
順序が変わるのが同一周期の完全直交な組だけであることから従う。

この還元は、数値的反射だけから新しい線束や very strongness を作るものではない。
特に式\eqref{eq:weyl-cross-block-move}の実現は、一般の非可換点では残された課題である。
可換点の許容形式変形では、通常の各標識の持上げによりこの実現が得られる。

% !TeX root = ../nc-cubic-unmark.tex
\section{可換三次曲面の構造対象とPicard作用}
\begin{theorem}[可換基準模型]\label{thm:commutative-structure}
滑らかな三次曲面$X$について、$\operatorname{coh}X$の許容構造対象は
ちょうど線束である。従って固定圏上の選択集合は$\Pic(X)\simeq\mathbb Z^7$である。
自然同型で割ったexact自己同値群は
\[
 G_X=\Pic(X)\rtimes\Aut(X)
\]
であり、構造対象の同型類に推移的に作用する。
従って $\Pic(X)/G_X$ は一点である。固定圏内の線束選択 $\Pic(X)$ と、
この軌道商を区別する。
\end{theorem}
\begin{proof}
可換滑らかな曲面で式\eqref{eq:gk-rank}を満たす層$E$は通常のrank-one
 torsion-free層である。二重双対$L=E^{**}$は線束で、
$E=L\otimes I_Z$、$n=\operatorname{length}Z\ge0$と書ける。
Riemann--Rochと$\operatorname{ch}E=e^{c_1(L)}(1-n[\mathrm{pt}])$から
\[
 \chi(E,E)=\chi(\OO_X)-2n=1-2n.
\]
例外性は$\chi(E,E)=1$を与えるので$n=0$、従って$E=L$である。
逆に線束$L$の自己Extは$H^*(X,\OO_X)$であり、$X$がrational del Pezzoなので例外的である。
$H=(-)\otimes\OO_X(-K_X)$、$-K_X$は豊富で、
Kodaira消滅とRiemann--Rochから式\eqref{eq:weak-hilbert}が成り立つ。
線束捩れは通常の豊富列を保つからAZ豊富性も成立する。

Bondal--Orlov\cite[Theorem 3.1]{BO}は反標準豊富な$X$の導来自己同値を
幾何学的自己同型、線束捩れ、次数shiftで記述する。
$t$-exact性は非零shiftを除き、上の$G_X$を得る。
線束捩れの推移性から、線束の同型類集合の軌道商は一点となる。
\end{proof}

\begin{proposition}[共通線束捩れはrelation表示に見えない]\label{prop:common-tensor}
exact同値$F$と有限thread$(E_i)$について、頂点冪等元と積を保つ同型
\[
 \End\Bigl(\bigoplus_iE_i\Bigr)
 \xrightarrow{\sim}\End\Bigl(\bigoplus_iFE_i\Bigr)
\]
がある。特に可換曲面上で全$E_i$を同じ線束$L$で捩っても、
矢の基底を輸送すればAOUの関係式点は変わらない。
先頭同定$E_0\simeq\OO_X$を固定した正規化には、一般の共通線束捩れは作用しない。
\end{proposition}
\begin{proof}
各Homへの$F$の作用は合成を保つ完全忠実な同型であり、直和すると主張を得る。
$F=(-)\otimes L$なら$E_0=\OO_X$は$L$へ移るため、
先頭同定を保つには$L\simeq\OO_X$が必要である。
\end{proof}
従ってAOUのrelation schemeをさらに一律に$\Pic(X)$で割る操作は、
この共通捩れが既に忘れられている点を考慮せずには正当化できない。
非可換圏では一般の対象同士をテンソルする操作そのものがなく、
自己同値による同一視は定理\ref{thm:forgetful-fibres}の集合論的軌道として別に扱う。

\section{Artin変形における対象の持上げと有限Weyl商}\label{sec:weyl}
以下の形式比較は、特殊ファイバーとの同定を保持する局所Artin底上で行う。
$\widehat{\Cw}$ は許容弱標識付き変形の同型類集合を取る関手を表す。
対象持上げにはInd-completionでのflat Abel変形の枠組みを使い、
相対perfect対象のRHom底変更、exactな$H_R$およびその底変更を仮定する。
これらを満たす変形を本節の\emph{許容相対変形}と呼ぶ。
全ての抽象dg変形にこれらのheart条件が自動的に備わるとはしない。

\begin{theorem}[構造対象の形式的な選択]\label{thm:formal-structure}
滑らかな三次曲面$X$の許容相対変形$\mathcal A_R$と$L\in\Pic(X)$に対し、
$L$のflatな持上げ$L_R$が存在し、その同型類は一意である。
$L_R$は相対許容構造対象であり、自己同型群は$R^\times$である。
従って同型類を取った集合値の形式変形関手では、構造対象の固定を忘れる写像は全単射であり、
固定圏上で全ての特殊ファイバーの選択を許せば枝は$\Pic(X)$で添字付けられる。
この主張は、任意の二つの枝を$\mathcal A_R$の自己同値が結ぶことを意味しない。
\end{theorem}
\begin{proof}
小拡大の核$I$に対する障害は$\Ext^2(L,L)\otimes I$、
同型類の自由度は$\Ext^1(L,L)\otimes I$である\cite[Theorem A and flat dual]{Lowen}。
両者が零なのでArtin長に関する帰納法で存在と一意性を得る。
相対RHomはperfectで、その特殊ファイバーは$\kk$が次数0にあるだけだから
$\RHom_R(L_R,L_R)\simeq R$となる。特殊ファイバーの同定を固定した自己同型は
$1+\mathfrak m_R$、固定しないものは$R^\times$である。
同様に$\RHom(L_R,H_R^mL_R)$の特殊ファイバーを計算すれば、
$m\ge0$の次数0集中と所定rank、$m<0$の$H^0$消滅が持ち上がる。

豊富性の確認も必要である。任意のnoetherian対象$M$は
$\mathfrak m_R$による有限filtrationを持ち、その商は特殊ファイバーの対象である。
flatな$H_R^{-n}L_R$に対する導来底変更と古典的Serre消滅から、
十分大きな$n$でこの各商への高次Extが零になる。
有限filtrationの長完全列により$M$についても同じ消滅が成立する。
$M/\mathfrak m_RM$への古典的生成全射を持ち上げると、
$\mathfrak m_RM$への$\Ext^1$消滅によって$M$への射が得られ、
余核はNakayamaにより零になる。これはAZ生成条件を与える。
任意の固定した全射の核に同じ$\Ext^1$消滅を使えばHom全射条件も得る。
rank条件は特殊ファイバー上で確認済みである。
\end{proof}

\begin{theorem}[通常の標識に限った形式Weyl商]\label{thm:formal-weyl}
$\Aut(X)=1$なる滑らかな三次曲面$X$を取り、通常の幾何学的標識の集合を
$\operatorname{Mark}(X)$とする。各$m$の標準九対象line-bundle helixで先頭を
$\OO_X$に正規化し、対象の同型類を取り、矢の基底と頂点スカラーの選択は数えない。
ここで形式関手は同型類の集合を取る。許容相対変形にそのhelixの持上げを加える形式関手を$\widehat{\mathfrak M}_m$とすると
\[
 \widehat{\mathfrak M}_m\simeq\widehat{\Cw}_{(\operatorname{coh}X,\OO_X)},\qquad
 \left(\coprod_{m\in\operatorname{Mark}(X)}\widehat{\mathfrak M}_m\right)/W(E_6)
 \simeq\widehat{\Cw}_{(\operatorname{coh}X,\OO_X)}.
\]
AOUの関係式の形式変形がこれらの許容相対条件を満たす範囲では、
その関係式表示が左辺を与える。Hochschild変形との比較が成立する形式関手の
miniversal底は$\operatorname{Spf}\kk[[q_1,\ldots,q_8]]$である。
\end{theorem}
\begin{proof}
標準九対象はline bundleの完全強例外列で、先頭は$\OO_X$である
\cite[(4.1), (9.44)]{AOU}。各対象は定理\ref{thm:formal-structure}と同じ障害消滅で
一意に持ち上がる。相対Homのperfect底変更から強例外性とHomのflat性を得る。
標準helixは$\mathbb Z$-生成geometric helixである。
De Deken--Lowen\cite[Theorems 4.5, 5.13, 5.15]{DDL}をInd-completionで適用すると、
helixの生成性と持上げ、Abel圏・helix代数・有限threadの変形の比較を得る。
ここでは同定した$H_R$がheartを保つことを許容相対条件に含めているので、
単なる導来同値からheartの同一性を推論してはいない。
反標準部分列の豊富性と切断数は前定理の証明で保持される。

$\operatorname{Mark}(X)$は古典的$W(E_6)$-torsorである\cite[\S2.3]{AOU}。
各chartを同じ忘却先と同定することで標識変更作用を定義でき、
持上げの一意性により合成は整合的である。
$\HH^1=\HH^3=0$、$\dim\HH^2=8$は命題\ref{prop:hh}による。
通常のHochschild変形理論との比較を満たす範囲で無障害性と8次元の形式底を得る。
\end{proof}

\begin{corollary}[固定相対 heart 上の標準型線束標識の網羅性]\label{cor:formal-all-line-markings}
特殊ファイバーを $\operatorname{coh}X$ と同定した許容相対変形 $\mathcal A_R$ と、
$L\in\Pic(X)$ の一意な持上げ $L_R$ を固定する。
定理\ref{thm:formal-weyl}の証明で用いる相対 perfect 底変更と helix の持上げ・生成の
条件を満たす範囲で、次の標識を数える：九対象は flat な相対線束であり、
特殊ファイバーの九対象は rank-one torsion-free で、AOU の三ブロック Euler 型を持ち、
先頭は $L_R$ である。対象ごとの同型で同一視し、圏の自己同値では割らない。
この集合は $X$ の通常の格子標識の集合と全単射であり、$W(E_6)$-torsor となる。
この固定圏上の主張には $\Aut(X)=1$ を仮定しない。
\end{corollary}
\begin{proof}
各通常標識の標準九線束を $L$ で共通 tensor し、前定理の証明によって
一意に持ち上げると、この集合の元を得る。
逆に任意の相対標識の特殊ファイバーを取る。
各対象は rank-one torsion-free で、特殊ファイバーの Euler 行列の対角成分は $1$ である。
定理\ref{thm:commutative-structure}の証明における
$\chi(E,E)=1-2\operatorname{length}Z$ より、各対象は通常の線束である。
補題\ref{lem:weyl-frame-reconstruction}によって一意な格子標識を復元でき、
可換曲面での線束の一意性から九対象はその標準九線束の共通 $L$ 捩れと一致する。
各対象の持上げの一意性により元の相対標識もこの標準持上げと一致する。
固定圏の自己同値を商に取っていないため、$\Aut(X)$ の作用はこの個数を減らさない。
\end{proof}
この系は特殊ファイバーの rank 条件を保つ上記相対標識の網羅性であり、
任意の OU 型 helix や任意の dg 変形を扱うものではない。

\begin{proposition}[有限 Weyl 商の集合論的な十分条件]\label{prop:global-weyl-criterion}
標識付き対象の集合 $M$ と弱標識付き対象の集合 $W$ の間の忘却 $p:M\to W$ を考える。
有限群 $W(E_6)$ が $M$ に作用し、$p$ が全射であり、
二点が同じ像を持つことと同じ $W(E_6)$ 軌道に属すことが同値なら、
$M/W(E_6)\to W$ は全単射である。
\end{proposition}
\begin{proof}
全射性は $p$ の全射性、単射性は仮定した軌道による同値関係そのものである。
\end{proof}
非可換の対象範囲でこの仮定を確認することが第一段階の課題である。
可換標識の形式持上げだけから、全ての helix が同一の有限 Weyl 軌道にあるとは結論しない。


% !TeX root = ../nc-cubic-unmark.tex
\section{Picard作用は固定した非可換fiberを保つとは限らない}\label{sec:picard-first-order}
対象の持上げと自己同値の持上げを区別するため、一次変形を直接計算する。
一般のFourier--Mukai同値はHochschild変形パラメータも移す
\cite[Theorem 1.1]{Toda}。以下はそのline-bundleの場合をČech表示で確認するものである。

\begin{proposition}[一次の標識変更障害]\label{prop:picard-first-order}
滑らかな三次曲面$X$の一次Abel変形を、HKRの規約
\[
 HH^2(X)=H^1(X,T_X)\oplus H^0(X,\wedge^2T_X),\qquad u=(\xi,\pi)
\]
で記述する。局所積は$a*b=ab+\frac\epsilon2\pi(da,db)$とする。
$L$を持ち上げた右加群の局所自己準同型環を使う標識変更は、変形類を
\begin{equation}\label{eq:picard-shear}
 r_L(\xi,\pi)=(\xi+\delta_\pi(L),\pi),\qquad
 \delta_\pi(L)=H^1(\pi^\sharp)(c_1(L))\in H^1(X,T_X)
\end{equation}
へ移す。ここで$\pi^\sharp(\alpha)(a)=\pi(\alpha,da)$、
$c_1(L)$は$d\log$のČech類である。
逆向きのMorita同値の規約では符号が反転する。
特に$(-)\otimes L$が\emph{同じ}変形の自己同値に持ち上がるためには
$\delta_\pi(L)=0$が必要である。
\end{proposition}
\begin{proof}
アフィン被覆上で$L$の遷移を$g_{ij}$と書く。
右加群の局所基底を$e_j=e_i*g_{ij}$とすれば、その自己準同型環の遷移は
$b\mapsto g_{ij}*b*g_{ij}^{-1}$である。
一次まで展開して
\[
 g*b*g^{-1}=b+\epsilon\,\pi(d\log g,db)
\]
を得る。元の環の遷移$1+\epsilon\xi_{ij}$と合成すれば
$\xi_{ij}+\pi^\sharp(d\log g_{ij})$となり、式\eqref{eq:picard-shear}を得る。
$L$の持上げに必要な一次の遷移補正は内的共役を一次では変えない。
三重交叉上のスカラー障害とその変更は$H^2(\OO_X)$に入るが、三次曲面では零である。
従ってこの計算はAbel変形の同型類に依存する。
固定変形の自己同値が持ち上がれば、標識変更前後の変形類は同じでなければならず、
$H^1(T_X)$成分の差は零となる。ここではこの必要条件から高次Artin底上の
十分性を結論していない。
\end{proof}

\begin{theorem}[三次曲面上での非自明性と階数]\label{thm:picard-obstruction-rank}
$\pi$の零因子$C\in|-K_X|$が滑らかなとき、複素線型写像
\[
 H^1(\Omega_X^1)\xrightarrow{H^1(\pi^\sharp)}H^1(T_X)
\]
の階数は3または4である。特に$\delta_\pi(L)\ne0$となる線束$L$が存在する。
階数4となる非空開部分があり、一例は
\begin{equation}\label{eq:picard-witness}
 X=V(x^3+y^3+z^3+txy+t^3),\qquad C=X\cap\{t=0\}.
\end{equation}
従って、全ての線束が対象として一意に持ち上がる場合でも、
それらのテンソル自己同値が同じ一次非可換変形に全て持ち上がるわけではない。
\end{theorem}
\begin{proof}
$\wedge^2T_X\simeq\OO_X(-K_X)$と$T_X\otimes K_X\simeq\Omega_X^1$から
\[
 0\longrightarrow\Omega_X^1\xrightarrow{\pi^\sharp}T_X
 \longrightarrow T_X|_C\longrightarrow0
\]
が得られる。$H^0(T_X)=0$、$h^1(\Omega_X^1)=7$なので、求める階数は
$7-h^0(T_X|_C)$である。楕円曲線$C$の正常完全列は
\[
 0\longrightarrow\OO_C\longrightarrow T_X|_C
 \longrightarrow\OO_C(C)\longrightarrow0.
\]
$\deg\OO_C(C)=3$より$h^0\OO_C(C)=3$、$h^1\OO_C(C)=0$である。
その接続写像$\partial:H^0(\OO_C(C))\to H^1(\OO_C)$の階数を$r\in\{0,1\}$とすると
$h^0(T_X|_C)=4-r$、従って求める階数は$3+r$となる。
$X$はrationalなので$c_1:\Pic(X)\otimes\mathbb C\simeq H^1(\Omega_X^1)$であり、
正の階数は少なくとも一つの線束で非零障害を与える。

式\eqref{eq:picard-witness}の$X$は滑らかである。
勾配が零なら$z=0$で、$t=0$では$x=y=0$となる。
$t\ne0$で$t=1$とすると$3x^2+y=3y^2+x=xy+3=0$となるが、
最初の二式の積は$9(xy)^2=xy$を与え、$xy=-3$と矛盾する。
$C$はFermat平面三次曲線である。
平面を$t=\epsilon z$へ変形すると、その平面内の曲線は
$x^3+y^3+z^3+\epsilon xyz=0$となる。
$xyz$はJacobian環$\mathbb C[x,y,z]/(x^2,y^2,z^2)$の次数3において非零なので、
この平面三次の一次変形は射影座標変更で自明化できない。
楕円曲線の次数3線束の変形は平行移動によって吸収されるため、
これは抽象曲線の非零Kodaira--Spencer類でもある。
従って$\partial$は全射、$r=1$であり階数は4である。
行列の階数4条件は開であるから非空開部分を得る。
この有限多項式の確認は\texttt{verify\_marking\_lattice.py}でも行う。
\end{proof}

\begin{remark}[ここから結論していないこと]
階数は$\Pic(X)\otimes\mathbb C$上の階数であり、
整数群$\ker(\delta_\pi:\Pic(X)\to H^1(T_X))$のrankをこの階数計算だけでは直ちに決定しない。非常に一般の組については定理\ref{thm:very-general-picard-kernel}で決定する。
またこの障害は一次であり、全次数での自己同値の持上げや大域Picard商の存在は未解決である。
一方、固定した一次変形上で全Picard群が自己同値として作用するという
素朴な予想を否定するには十分である。
特に$\Aut(X)=1$の許容相対変形では、$\delta_\pi(L)\ne0$なら
$\OO_R$と$L_R$を移す自己同値は存在しない。
なぜならその特殊化は定理\ref{thm:commutative-structure}により
$(-)\otimes L$でなければならないからである。
\end{remark}


\begin{proposition}[反標準線束の障害は零]\label{prop:canonical-kernel}
$C=(\pi=0)\in|-K_X|$が滑らかなとき
\[
 \delta_\pi(\OO_X(-K_X))=0.
\]
さらに$Q=K_X^\perp\subset\Pic(X)$とすると
\[
 \operatorname{rank}_{\mathbb C}
 \bigl(\delta_\pi|_{Q\otimes\mathbb C}\bigr)
 =\operatorname{rank}_{\mathbb C}\delta_\pi
 \in\{3,4\}.
\]
従って一般の$\pi$では、$E_6$ root translationの中に
固定一次変形の自己同値としては持ち上がらないものがある。
\end{proposition}
\begin{proof}
$C$の局所方程式を$s_i$とし、$g_{ij}=s_i/s_j$を
$\OO_X(C)$の遷移の一つの規約とする。
$\pi$は局所的に$s_i$で割り切れるので
$v_i=\pi^\sharp(d\log s_i)$は各開集合上で正則なベクトル場である。
そのČech差は$\pi^\sharp(d\log g_{ij})$だから、
式\eqref{eq:picard-shear}の$\delta_\pi(\OO_X(C))$は零となる。
$\Pic(X)\otimes\mathbb C=(Q\otimes\mathbb C)\oplus\mathbb C K_X$
は$(-K_X)^2=3$から従う。
$K_X$方向で$\delta_\pi=0$なので、全Picardと$Q$への制限の階数は等しい。
階数は定理\ref{thm:picard-obstruction-rank}から従い、
非零の根格子方向は同じ一次変形における自己同値としては持ち上がらない。
\end{proof}


\begin{theorem}[非常に一般の一次Picard障害の整数核]\label{thm:very-general-picard-kernel}
滑らかな複素三次曲面$X$と非零のPoisson構造
$\pi\in H^0(X,-K_X)$の組$(X,\pi)$を動かす。
非常に一般の組に対して、群準同型
$\delta_\pi:\Pic(X)\to H^1(X,T_X)$の整数核は
\begin{equation}\label{eq:very-general-picard-kernel}
 \boxed{\ker\delta_\pi=\mathbb Z K_X}
\end{equation}
である。ここで「非常に一般」とは、適切な既約パラメータ空間内の
可算個の真の閉部分集合の合併の外を意味する。
複素線型写像$\delta_\pi\otimes\mathbb C$は一般には階数4であり、
その3次元核がPicard整数格子の階数3の部分格子を含むとは限らない。
\end{theorem}
\begin{proof}
六本の順序付きで互いに交わらない直線によって標識を付けた滑らかな三次曲面の
モジュライ$\mathcal M^{\rm m}$は、一般位置の六点配置の
$\PGL_3$商として既約であり、$W(E_6)$は標識変更によって作用する
\cite[\S\S2.5--2.6]{DGK}。
この上でPoisson構造$\pi$の零因子が滑らかな部分の非零全空間を
$\mathcal P^{\rm m}$とする。これは既約である。
標準Picard格子$\Lambda=I_{1,6}$は標識によって自明化され、
\[
 \Lambda_\mathbb C=\mathbb C K\oplus Q_\mathbb C,
 \qquad Q=K^\perp\simeq E_6
\]
となる。

$\mathcal P^{\rm m}$のすべての点$b=(X,\varphi,\pi)$についての
$\delta_b$の核の共通部分を
\[
 N=\bigcap_{b\in\mathcal P^{\rm m}}
       \ker(\delta_b:\Lambda_\mathbb C\to H^1(X,T_X))
\]
と置く。これは標識変更の自然性から$W(E_6)$不変である。
命題\ref{prop:canonical-kernel}により$\mathbb CK\subset N$である。
$W(E_6)$は$Q_\mathbb C$上で既約に作用する。
実際、任意の非零な不変部分空間$U$のベクトルがある根$\alpha$と
交叉積非零なら、反射$s_\alpha(v)-v=(v,\alpha)\alpha$から
$\alpha\in U$となる。根のWeyl軌道は$Q_\mathbb C$を張るので
$U=Q_\mathbb C$である。
定理\ref{thm:picard-obstruction-rank}の階数4の非空開部分があるため、
$N$は$\Lambda_\mathbb C$全体ではない。
従って$N=\mathbb CK$である。

各$D\in\Lambda\setminus\mathbb ZK$に対し、
$\delta_b(D)$は$\mathcal P^{\rm m}$上で恒等的に零ではない
（$K$は原始元である）。
これは代数的ベクトル束の切断なので、零点集合$Z_D$は真の閉部分集合である。
$\Lambda$は可算なので、
$\mathcal P^{\rm m}\setminus\bigcup_{D\notin\mathbb ZK}Z_D$
上では非反標準の全整数クラスの障害が非零となる。
他方、$\mathbb ZK$は常に障害の核に入る。
これで式\eqref{eq:very-general-picard-kernel}を得る。
\end{proof}

\begin{corollary}[非常に一般の非可換fiberでの構造対象の同型類]\label{cor:first-order-pointings}
$X$は自己同型群が自明な滑らかな三次曲面とし、
$\mathcal A_\varepsilon$は特殊ファイバー$\operatorname{coh}X$との同定を持つ
許容相対一次Abel変形（$\varepsilon^2=0$）とする。
HKRのPoisson成分$\pi$を含む組$(X,\pi)$が
定理\ref{thm:very-general-picard-kernel}の非常に一般の条件を満たすとき、
exactな$\mathbb C[\varepsilon]/(\varepsilon^2)$線型自己同値を
特殊ファイバーへ制限する準同型
\[
 \pi_0\Aut_{\rm ex}(\mathcal A_\varepsilon)
       \longrightarrow\Pic(X)
\]
の像は$\mathbb Z[-K_X]$に一致する。
従って特殊ファイバーで線束となる許容構造対象の持上げを
固定圏内のexact自己同値で同一視した同型類の集合は
\begin{equation}\label{eq:very-general-pointings}
 \Pic(X)/\mathbb ZK_X
 \ \simeq\ Q(E_6)^\vee
 \ \simeq\mathbb Z^6
\end{equation}
である。これは同型類の集合の主張であり、
相対的なモジュライ構造についての主張ではない。
\end{corollary}
\begin{proof}
定理\ref{thm:formal-structure}により、各線束$L$は
構造対象$L_\varepsilon$へ一意に持ち上がる。
Bondal--Orlovによるexact自己同値の可換特殊化の分類
（定理\ref{thm:commutative-structure}）と$\Aut(X)=1$から、
任意の相対exact自己同値の特殊化は$(-)\otimes L$となる。
命題\ref{prop:picard-first-order}の必要条件は
$\delta_\pi(L)=0$なので、前定理により$L$の類は$\mathbb ZK_X$に属す。
逆に、相対Serre作用から定まる$t$-exact自己同値
$H_\varepsilon=\mathsf S_\varepsilon^{-1}[2]$は
特殊ファイバー上で$(-)\otimes\OO_X(-K_X)$となる。
従って$H_\varepsilon$の全冪が実際の自己同値を与え、
特殊化像は$\mathbb Z[-K_X]$と一致する。
線束持上げの一意性とこの作用から同型類の商が
$\Pic(X)/\mathbb ZK_X$となる。
最後の格子同型は命題\ref{prop:picard-root-index}による。
\end{proof}

\begin{remark}
この結果は「非常に一般」という条件と一次Artin底、
許容相対heart・Serreの仮定を持つ。
任意の有限型の非可換大域成分の一般点に直ちに適用するものではなく、
高次変形での自己同値の全持上げやAOUの標識付きGIT商の有限性も結論しない。
とくに$Q(E_6)^\vee$は整数格子の商であり、
Poisson一次変形の底に新しい連続的な6次元成分を作ったわけではない。
\end{remark}

\section{有限Weyl群、root latticeとPicard格子の区別}
\begin{proposition}[反標準方向を除いたPicard格子]\label{prop:picard-root-index}
$\Lambda=\Pic(X)=\mathbb Zh\oplus\bigoplus_{i=1}^6\mathbb Ze_i$に
$h^2=1$、$e_i^2=-1$、他の積0を入れ、$K=-3h+\sum_i e_i$とする。
$Q=K^\perp$は$E_6$ root latticeであり、
\[
 0\longrightarrow Q\longrightarrow\Lambda
 \xrightarrow{D\mapsto(-K)D}\mathbb Z\longrightarrow0,
 \qquad \Lambda/(Q+\mathbb ZK)\simeq\mathbb Z/3.
\]
交叉積による制限写像は$\Lambda/\mathbb ZK\simeq Q^\vee$を与える。
従って反標準tensor方向を同一視した後でも、root translationと
Picard translationの間には指数3の相違がある。
\end{proposition}
\begin{proof}
次数写像は$e_i$を1へ送るので全射で、核は定義により$Q$である。
$K$の次数は$-3$なので商は$\mathbb Z/3$となる。
$Q$は原始部分格子、$\Lambda$はunimodularだから、
任意の整数線型形式$Q\to\mathbb Z$は$\Lambda$へ延長でき、
交叉積により$\Lambda$の元で表される。
制限写像$\Lambda\to Q^\vee$の核は$\mathbb ZK$である。
単純根$e_1-e_2,\ldots,e_5-e_6,h-e_1-e_2-e_3$のGram行列は
負の$E_6$ Cartan行列で行列式3を持つことからも指数を確認できる。
\end{proof}

\cite{OWR}で考えるaffine Weyl群のtranslation部分は$Q$であり、
全$\Pic(X)$と同じではない。
$\Lambda/\mathbb ZK$をtranslationとする記述はweight lattice型の拡大となる。
ただし命題\ref{prop:common-tensor}により共通tensorはrelation表示では既に見えない。
従ってこの指数3だけからモジュライの三重被覆や三つの成分を主張してはならない。
さらに定理\ref{thm:very-general-picard-kernel}が示すように、非常に一般の一次変形では整数Picard群の障害の核は反標準方向だけである。非可換変形上のPicard作用は
一般には底の変形パラメータをも動かす。全線束の存在と、固定圏の自己同値による軌道は分けて調べる。

\section{一次 Picard shear の代数的不変式と粗商の障害}
\label{sec:picard-invariant-quotient}

この節では線形一次変形空間の商を厳密に計算する。これは AOU の
GIT 商そのものを計算するものではない。滑らかな三次曲面 $X$ に対し
\[
 V:=H^1(X,T_X),\qquad P:=H^0(X,-K_X),\qquad
 E:=V\oplus P=\HH^2(X)
\]
と書く。どちらも四次元である（命題\ref{prop:hh}）。
$\Lambda=\Pic(X)$ とおき、$L\in\Lambda$ に対して
\begin{equation}\label{eq:picard-linear-shear}
 r_L(\xi,\pi):=(\xi+\delta_\pi(L),\pi),\qquad
 \delta_\pi(L):=H^1(\pi^\sharp)c_1(L)
\end{equation}
と定める。$\delta_\pi(L)$ は $\pi$ と $L$ の第一 Chern 類に
それぞれ線形なので、$r_L$ は $E$ の線形自己同型であり
$r_Lr_M=r_{L\otimes M}$ となる。この表示は
命題\ref{prop:picard-first-order}で確認した Morita 標識変更の一次作用である。

\begin{lemma}[加法群における格子平行移動の Zariski 閉包]
\label{lem:translation-closure}
特性零の体 $k$ 上の有限次元線形空間 $W$ と有限個の元
$v_1,\ldots,v_s\in W$ に対し、$\sum_i\mathbb Z v_i$ の
$W$ の加法群における Zariski 閉包は
$\operatorname{span}_k\{v_1,\ldots,v_s\}$ に等しい。
\end{lemma}
\begin{proof}
$\Gamma=\sum_i\mathbb Z v_i$ とする。$\Gamma$ 上で消える
多項式 $f$ を取り、一つの $\gamma\in\Gamma$ を固定する。
$\gamma+n v_i\in\Gamma$ より $f(\gamma+t v_i)$ は
整数 $t=n$ のすべてで零となるので恒等的に零である。
この操作を各 $i$ について繰り返すと
$\gamma+\sum_i t_i v_i$ 上で $f$ は恒等的に零である。
特に $\gamma=0$ と置けば $f$ は張る部分空間上で零になる。
逆向きの包含は自明である。
\end{proof}

\begin{theorem}[Picard shear の多項式・有理不変式]
\label{thm:picard-linear-invariants}
$X$ は滑らかな複素三次曲面であり、$\pi\in P$ が滑らかな零因子を持ち、
$H^1(\pi^\sharp)$ の階数が $4$ になるものが存在すると仮定する。
$\Lambda=\Pic(X)$ は式\eqref{eq:picard-linear-shear}によって
$E\simeq\mathbb A^8$ に作用し、この作用の不変式は
\begin{equation}\label{eq:picard-invariants}
 \boxed{\ \mathbb C[E]^\Lambda=\mathbb C[P],\qquad
 \mathbb C(E)^\Lambda=\mathbb C(P)\ }.
\end{equation}
実際、階数が $4$ となる任意の $\pi$ に対し、$\Lambda$ の軌道は
各 $V\times\{\pi\}$ の中で Zariski 稠密である。
特に、$\pi=0$ 上では全 $r_L$ が恒等であるにもかかわらず、
\emph{一次作用の affine な不変式商} $E\to\Spec\mathbb C[E]^\Lambda\cong P$
は $V\times\{0\}$ 全体を一点に写す。
\end{theorem}
\begin{proof}
$H^1(\pi^\sharp)$ の階数 $4$ となる $\pi$ の集合 $P^{\rm surj}$ は
非空 Zariski 開である。$c_1(\Pic X)$ は
$H^1(X,\Omega_X^1)$ を複素線形に張るので、
各 $\pi\in P^{\rm surj}$ について
$\delta_\pi(\Lambda)$ の複素線形包は $V$ 全体になる。
補題\ref{lem:translation-closure}により、この加法部分群は $V$ で
Zariski 稠密である。$f\in\mathbb C[E]^\Lambda$ を取ると、
$\pi\in P^{\rm surj}$ を固定した $f(\cdot,\pi)$ は
すべての $\delta_\pi(\Lambda)$ による平行移動で不変だから
$V$ 上の定数である。$P^{\rm surj}$ が稠密であるので
$f$ は $\xi$ に依存せず、$f\in\mathbb C[P]$ である。
逆向きの包含は $r_L$ が $\pi$ を固定することから従う。

有理不変式 $f\in\mathbb C(E)^\Lambda$ についても、各
$\pi\in P^{\rm surj}$ の一般点で $f(\cdot,\pi)$ を評価すれば
同じ結論を得る。厳密には $\mathbb C(P)$ 上で $f$ を
$V$ の有理関数とみなし、各 $L$ による整数平行移動の不変性から
方向 $\delta_\pi(L)$ に沿うすべての有理微分が零であることを得る。
$\delta_\pi(\Lambda)$ は一般の $\pi$ で $V$ を張るため
$V$ に関するすべての偏微分が零となり、
特性零により $f\in\mathbb C(P)$ が従う。
\end{proof}

\begin{corollary}[分離的な一次粗商の限界]
\label{cor:separated-linear-quotient}
上の定理の条件の下で、$Y$ を複素数体上の有限型分離代数空間とする。
Picard shear に不変な射 $F:E\to Y$ は
\[
 E\xrightarrow{\operatorname{pr}_P}P\xrightarrow{g}Y
\]
と一意に因子化する。従って $F$ は可換方向 $V\times\{0\}$ の
異なる点を区別できない。
\end{corollary}
\begin{proof}
$\pi\in P^{\rm surj}(\mathbb C)$ と $\xi\in V$ を取る。
$\Lambda$ 軌道上で $F$ は定数であり、その軌道は
$V\times\{\pi\}$ に Zariski 稠密である。
$Y$ は有限型なので複素点は閉点であり、射の連続性から
$F$ は $V\times\{\pi\}$ 全体で定数となる。
零切断 $s:P\to E$ に対し $F$ と
$F\circ s\circ\operatorname{pr}_P$ は
稠密開 $V\times P^{\rm surj}$ で一致する。
$Y$ は分離なので両者の等化子は閉集合であり、
$E$ 全体で一致する。$g=F\circ s$ とすればよい。
\end{proof}

\begin{remark}[適用範囲と AOU への意味]
\label{rem:picard-linear-scope}
定理\ref{thm:picard-linear-invariants}と系
\ref{cor:separated-linear-quotient}は、$X$ の\emph{一次 Hochschild
変形空間}上の線形化した Picard 作用についての厳密な結論であり、
AOU の GIT スキーム上で同じ作用が全域に定義されることや、
非可換 Abel 圏の大域粗モジュライが存在しないことを意味しない。
特に一般の $\pi$ に対する整数核 $\mathbb ZK_X$
（定理\ref{thm:very-general-picard-kernel}）は軌道の\emph{集合論的}
分類を与える一方、ここでの Zariski 稠密性は
有限型分離な粗商の\emph{情報喪失}を与える。
この相違が、AOU 関係式の通常の GIT 商だけから
無標識 Abel 圏のモジュライを取り出す際の主要な障害となる。
\end{remark}

\section{現在の課題}
最優先は二点である。第一は系\ref{cor:weyl-one-move}の三対象交換を
固定反標準 heart 内で構成し、線束性・完全性・very strongness を示すことである。
数値上一意な候補が分かっても、核・余核・複体の cohomology が
必要な対象になることは別に確認する。
第二は全数値的線束類の実現と weak 許容性を決定することである。
零次数 root 型単純対象の存在・排除、数値的点に対する Ext 消滅から
一般の生成へ進む条件、および正の反標準自己 twist の消滅を調べる。
AS heart との比較はこの二問題に必要な範囲だけで行う。
自己同値軌道、相対族、群oid・stack はこの分類の後の段階とする。

\appendix
\section{可換三次曲面のHochschild変形}\label{sec:hh}
\begin{proposition}\label{prop:hh}
滑らかな可換三次曲面$X$について
\[
 \HH^1(X)=0,\qquad \dim\HH^2(X)=4+4=8,\qquad \HH^3(X)=0.
\]
従ってHochschild複体で制御される通常の標識付き形式変形は無障害で、
そのminiversalな底は8次元形式円板である。
\end{proposition}
\begin{proof}
HKRにより$\HH^2=H^1(T_X)\oplus H^0(-K_X)\oplus H^2(\OO_X)$。
$H^2(\OO_X)=0$、$h^0(-K_X)=4$である。
接束の正常完全列とEuler完全列から$H^0(T_X)=0$および
$h^1(T_X)=h^0(\OO_X(3))-15=19-15=4$を得る。
前者の消滅は、$F$の偏導関数が二次正則列であり、一次係数のsyzygyがないことにもよる。
$H^2(T_X)=0$は正常完全列の高次コホモロジー、
$H^1(-K_X)=0$は$X$の直線束の中間コホモロジー消滅から従う。
$\wedge^3T_X=0$なので$\HH^3=0$。
同様に$\HH^1=H^0(T_X)\oplus H^1(\OO_X)=0$。
通常のHochschild変形理論では障害は$\HH^3$に入るため無障害性が従う。
この計算は反標準環の全域AS性や量子化の代数化を証明するものではない。
\end{proof}



\section{全関係式空間上の普遍代数と形式変形の比較}
\label{sec:aou-universal}
本節では、AOU の有限 quiver の特殊な三層構造を用い、
関係式の変形が有限次元代数の変形をどこまで尽くすかを決定する。
これは\emph{有限次元代数・導来圏}についての結果であり、
非可換三次曲面の Abel heart の全 GIT 境界への延長を意味しない。

\subsection{三層 quiver の普遍代数}
頂点集合を $Q_0=U\sqcup W\sqcup Z$ とし、各層に三頂点を置く。
$U$ から $W$、$W$ から $Z$ へは各組にちょうど一本の矢があり、
それ以外の矢はない。ここでは道の積を左から右へ読み、
右射影加群 $e_vB$ が頂点 $v$ から出る道で張られる規約を採用する。
これは AOU の $Q$ を矢の始点・終点で
並べ直した表示である\cite[\S4.2, Figure~4.1]{AOU}。
各 $(u,z)\in U\times Z$ について、長さ二の道の空間
$P_{uz}:=e_u(kQ)e_z$ は三次元である。
\[
 R_Q=\prod_{(u,z)\in U\times Z}\mathbb P(P_{uz})
       \simeq(\mathbb P^2)^9
\]
上の九つの tautological line を $\mathscr I_{uz}\subset
P_{uz}\otimes\mathcal O_{R_Q}$ とする。
それらが生成する両側イデアル $\mathscr I$ により
\[
 \mathscr B=(\mathcal O_{R_Q}Q)/\mathscr I
\]
を定める。AOU の積の規約と異なる場合は反対代数を取ればよく、
以下の変形・大域次元・Hochschild cohomology の結論は変わらない。

\begin{theorem}[全関係式空間上の有限代数と導来圏]
\label{thm:universal-directed-family}
$\mathscr B$ は $R_Q$ 上階数 $45$ の局所自由な代数である。
すべての幾何学的点 $r\in R_Q$ について $B_r=\mathscr B\otimes k(r)$ は
大域次元がちょうど $2$ の有限次元代数であり、
$\operatorname{Perf}(B_r)$ は smooth proper な dg 圏となる。
この族は矢の基底変換群 $G_Q=(\mathbb G_m)^{18}$ に同変である。
従って同じ矢の基底変更軌道にある関係式は、同型な有限代数を与える。
特にこの結論は AOU の GIT 半安定条件にも AS 正則性にも依存しない。
\end{theorem}
\begin{proof}
$Q$ に長さ三以上の道はなく、すべての関係式は長さ二である。
従って $\mathscr B$ の頂点・矢の成分はそれぞれ階数 $9,18$ の自由加群であり、
各 $(u,z)$ の長さ二の成分は $P_{uz}/\mathscr I_{uz}$ という
階数二の局所自由加群である。よって階数は $9+18+9\cdot2=45$ である。

固定した幾何学的ファイバーで、$P_v=e_vB_r$ を右射影加群、
$S_v$ をその単純商とする。$z\in Z$ に対し $S_z=P_z$ は射影的である。
$w\in W$ についてはその radical が三つの $P_z$ の直和なので
$S_w$ は長さ一の射影解消を持つ。
$u\in U$ に対し、$\operatorname{rad}P_u$ は三つの $P_w$ の像で生成され、
その間の核は各 $z\in Z$ に一つずつ与えられた非零関係式で生成される。
各関係式の終点は異なり、$P_z=S_z$ なので、この核は
$\bigoplus_{z\in Z}P_z$ に同型である。
従って
\[
 0\longrightarrow\bigoplus_{z\in Z}P_z
 \longrightarrow\bigoplus_{w\in W}P_w
 \longrightarrow P_u\longrightarrow S_u\longrightarrow0
\]
は射影解消となる。すべての単純加群の射影次元は高々二だから
$\operatorname{gldim}B_r\le2$ である。非零の二次関係式はこの射影解消に非零の第二syzygyを与え、
写像は radical 内に入るので解消は極小である。従って
$\operatorname{gldim}B_r=2$ である。
有限次元・有限大域次元の代数は homologically smooth で proper なので
その perfect dg 圏も smooth proper となる。
相対的にも、頂点の分離代数 $S=\mathcal O_{R_Q}^{Q_0}$、
矢の $S$-双加群 $V$、九つの普遍関係式の $S$-双加群 $J$ を使う標準の
$\mathscr B$-双加群複体
\[
 0\longrightarrow \mathscr B\otimes_S J\otimes_S\mathscr B
 \longrightarrow \mathscr B\otimes_S V\otimes_S\mathscr B
 \longrightarrow \mathscr B\otimes_S\mathscr B
 \longrightarrow \mathscr B\longrightarrow0
\]
がある。長さ三以上の道がないため関係式の重なりによる高次項はなく、
各幾何学的ファイバーでこの複体は完全となる。
完全性は長さ二までの道だけで直接確認できる。各ファイバーで右端の積写像は全射で、
その核は矢から生成されるため二番目の写像の像に等しい。
$\dim(B\otimes_S B)=108$、$\dim(B\otimes_S V\otimes_S B)=72$、
$\dim(B\otimes_S J\otimes_S B)=9$ だから、
二番目の写像の核は九次元である。
関係式を一つの道の二箇所で切って得る双加群写像
$J\to B\otimes_S V\otimes_S B$ は、各関係式が非零で、
九つの始点終点が異なることから単射である。
従って上の複体は各ファイバーで完全である。
各項は基底上で局所自由なので fiberwise exactness と Nakayama により
相対的にも完全である。ゆえに普遍対角双加群は perfect であり、
族は相対的にも smooth proper である。
矢の基底変換は普遍イデアルを輸送して $\mathscr B$ の同型を誘導するから、
族は $G_Q$-同変である。
\end{proof}

\begin{remark}[何がまだ Abel 圏ではないのか]
\label{rem:tilting-heart-difference}
可換三次曲面 $X$ の AOU 点では
$\operatorname{Perf}(B_r)\simeq D^b(\operatorname{coh}X)$ となるが、
$\operatorname{mod}B_r$ と $\operatorname{coh}X$ が Abel 圏として同値なのではない。
実際、前者は有限長の Abel 圏である一方、後者は有限長圏ではない。
従って全関係式空間に普遍的に存在する標準 heart
$\operatorname{mod}B_r$ を、そのまま求める非可換三次曲面の
無標識 Abel 圏と呼ぶことはできない。
必要なのは、導来圏上に\emph{三次曲面型の別の heart} を認識・構成し、
Serre 作用と helix に両立させることである。
\end{remark}

\subsection{境界上の Hochschild cohomology の跳び}
各関係式点 $r$ において、その $(u_i,z_j)$ 成分を
三つの長さ二の道の基底に関する斉次座標
$[c_{ij0}:c_{ij1}:c_{ij2}]$ と書き、
$S_{ij}(r)=\{m\mid c_{ijm}\ne0\}$ とする。
18 本の矢にそれぞれ変数 $a_{im}$（$u_i\to w_m$）と
$b_{mj}$（$w_m\to z_j$）を割り当てる。
整数係数行列 $D_r$ は、各 $(i,j)$ ごとに $m_0\in S_{ij}(r)$ を一つ固定し、
$m\in S_{ij}(r)\setminus\{m_0\}$ について
\[
 (a_{im}+b_{mj})-(a_{im_0}+b_{m_0j})
\]
を行とする行列である。この階数は $m_0$ の選択に依存しない。

\begin{theorem}[境界点での Hochschild 次元の組合せ公式]
\label{thm:aou-hh-support}
任意の $r\in R_Q(\mathbb C)$ に対し、$\rho(r)=\operatorname{rank}_{\mathbb C}D_r$ とすると
\[
\begin{aligned}
 HH^0(B_r)&\simeq\mathbb C,&\quad \dim HH^1(B_r)&=10-\rho(r),\\
 \dim HH^2(B_r)&=18-\rho(r),&\quad HH^n(B_r)&=0\quad(n\ge3).
\end{aligned}
\]
特に $\dim HH^2(B_r)-\dim HH^1(B_r)=8$ は全関係式空間で一定である。
九つの関係式の全係数が非零なら $\rho(r)=10$ であり
$(\dim HH^1,\dim HH^2)=(0,8)$ となる。
すべての関係式が一つの道のみからなる単項式関係式点では
$\rho(r)=0$ であり $(\dim HH^1,\dim HH^2)=(10,18)$ となる。
\end{theorem}
\begin{proof}
定理\ref{thm:universal-directed-family} の証明の相対双加群解消を
$r$ で特殊化し、$\operatorname{Hom}_{B_r^e}(-,B_r)$ を適用する。
得られる Hochschild 複体は
\[
 C^0=\bigoplus_{v\in Q_0}e_vB_re_v\simeq\mathbb C^9,
 \quad C^1=\bigoplus_{a\in Q_1}e_{s(a)}B_re_{t(a)}\simeq\mathbb C^{18},
 \quad C^2=\bigoplus_{(i,j)}e_{u_i}B_re_{z_j}\simeq\mathbb C^{18}.
\]
$C^n=0$ ($n\ge3$) である。
最初の微分 $C^0\to C^1$ は連結な quiver の頂点スカラーの差なので
階数は $9-1=8$ である。
次の微分は、矢をその係数倍だけ無限小に変化させた際の
九つの関係式の変化を、各 $P_{ij}/\langle r_{ij}\rangle$ に射影するものとなる。
$[c_{ij0}:c_{ij1}:c_{ij2}]$ の零でない成分については
$(a_{im}+b_{mj})$ が同じ値なら関係式は不変であり、
零の成分には条件がない。従ってこの微分の階数は
$\rho(r)$ に一致する。
以上から
$\dim HH^1=(18-\rho)-8$、$\dim HH^2=18-\rho$ を得る。

全係数非零の場合は、前節の三層構造による方程式
$a_{im}+b_{mj}=c_{ij}$ の核が頂点スカラーの八次元空間に
一致するので、$\rho=18-8=10$ である。
単項式関係式の場合は $S_{ij}$ がすべて一元集合なので
$D_r$ の行がなく、$\rho=0$ である。
\end{proof}

\begin{remark}[境界で跳ぶのは何か]
この公式は、関係式の無限小安定化群と有限代数の変形空間が
境界で同時に大きくなることを明示する。
しかし $HH^1(B_r)$ は $B_r$ の\emph{導来自己同値の無限小変形}を測るものであり、
ここから非可換三次曲面の Abel heart の自己同値群を無条件には復元できない。
\end{remark}

\subsection{Artin 変形を九つの関係式で尽くすこと}
\begin{theorem}[関係式の形式的完備性]
\label{thm:aou-relation-formal-completeness}
$k$ を代数閉体、$r\in R_Q(k)$、$B_0=B_r$ とする。
局所 Artin $k$-代数 $(R,\mathfrak m)$ 上の $R$-平坦な
結合代数 $B_R$ で、$B_R\otimes_Rk\cong B_0$ と同定したものを考える。
原始冪等元をラベル付きで持ち上げておけば、$B_R$ は
$R_Q$ の $r$ を通る $R$-点に付随する $R$-代数と同型である。
さらに、冪等元を固定し、特殊ファイバーで恒等となる同型による
$B_R$ の同値関係は、18本の矢の $R^\times$ による
基底変換（特殊ファイバーで恒等）で関係式を同一視する関係に一致する。
冪等元の持上げを取り替えた場合も、$1+\mathfrak mB_R$ による
内部共役で調整できる。従って Artin 局所上の
\emph{有限次元代数の変形}は九関係式の変形によって尽くされる。
\end{theorem}
\begin{proof}
冪等元の持上げは冪零イデアルに沿って存在する。
$e_vB_Re_w$ は $B_R$ の直和因子であり、有限自由 $R$-加群で、
その階数は特殊ファイバーの対応する成分の次元に等しい。
特に隣接層の $18$ 成分は階数一、$U$ から $Z$ への九成分は階数二、
その他は零または頂点の階数一である。

各隣接層に対応する階数一成分の基底を、特殊ファイバーの矢の基底から
持ち上げて選ぶ。これらは $B_R$ を生成する：特殊ファイバーで生成し、
$B_R$ は有限 $R$-加群なので Nakayama の補題が適用できる。
各 $(u,z)$ について長さ二の三つの道の積は
$R^3\twoheadrightarrow e_uB_Re_z\simeq R^2$ を与える。
特殊ファイバーで全射なので Artin 局所上でも全射であり、
その核は階数一の直和因子 $\mathscr I_{uz,R}$ である。
長さ三以上の道は存在しないので、
これら九つの核で生成されるイデアルによる剰余は $B_R$ に同型となる。
これは $R_Q(R)$ の点である。

冪等元を固定する代数同型は、それぞれの階数一の矢の成分上で
単元倍として作用する。それ以外の成分は矢の積で生成されるため、
同型は18本の矢の基底変更に一意に帰着する。
また冪零イデアルに沿う二組の完全直交冪等元の持上げは
$1+\mathfrak mB_R$ の元で共役である。
従って主張が従う。
\end{proof}

\begin{proposition}[可換点での八次元 Kodaira--Spencer 比較]
\label{prop:aou-kodaira-spencer-iso}
$X$ を滑らかな三次曲面、$r$ をその標準 AOU 九対象 thread の関係式点とし、
$B_0=\operatorname{End}(\bigoplus E_v)$ と置く。
関係式の全係数が非零である稠密トーラス上では、
矢の基底変更の Lie 代数が $T_rR_Q$ に作用する接写像は
階数 $10$ で、核は頂点基底変更の $8$ 次元空間である。
従って商の接空間は $18-10=8$ 次元であり、
その有限次元代数変形への Kodaira--Spencer 写像は
\[
 T_{[r]}M_Q^\circ\xrightarrow{\ \sim\ }HH^2(B_0)
 \xrightarrow{\ \sim\ }HH^2(X)
 =H^1(X,T_X)\oplus H^0(X,-K_X)
\]
という同型である。
ここで第一の同型は標識付き\emph{有限代数}の変形問題についてのものであり、
全 AOU 関係式点に相対的な三次曲面 Abel heart が存在するという主張ではない。
\end{proposition}
\begin{proof}
$U,W,Z$ の頂点をそれぞれ $u_i,w_m,z_j$ と記す。
最初の矢の無限小スカラーを $a_{im}$、次の矢のものを $b_{mj}$ とすると、
九つの各関係式の三つの係数への作用は
$\bigl(a_{im}+b_{mj}\bigr)_{m=0,1,2}$ の共通スカラーを除いたものになる。
核はすべての $i,j,m$ に対して
\[
 a_{im}+b_{mj}=c_{ij}\quad(m=0,1,2)
\]
を満たすものからなる。$m=0$ と比較すれば
$a_{im}-a_{i0}=b_{0j}-b_{mj}$ は $i,j$ に依らない。
この値を $d_m$ と書けば
$a_{im}=a_{i0}+d_m$、$b_{mj}=b_{0j}-d_m$ となる。
これは各頂点のスカラーによる基底変更に他ならず、
核の次元は $3+3+3-1=8$、作用の階数は $18-8=10$ である。
同じ整数係数行列について Smith 標準形が
$\operatorname{diag}(1^{\times10},0^{\times8})$ となることは
\path{computations/verify_aou_relation_tangent.py} で検算できる。

定理\ref{thm:aou-relation-formal-completeness}により、
関係式の一次変形を矢の基底変換で割った空間は
$B_0$ の結合代数一次変形を同型で割った空間、すなわち $HH^2(B_0)$ である。
九対象は完全強例外 collection であるから
$D^b(\operatorname{mod}B_0)\simeq D^b(\operatorname{coh}X)$。
Hochschild cohomology の導来 Morita 不変性と命題\ref{prop:hh}から
$HH^2(B_0)\simeq HH^2(X)$、その次元は八である。
これにより同型が従う。
\end{proof}

\begin{remark}[到達した比較と残る心の問題]
\label{rem:relation-vs-heart}
この命題は AOU 関係式の八つの一次方向が\emph{すべて}
$HH^2(X)$ を与えることを証明している。
他方、(i) その一般点での AS-regular $\{0,1,2\}\times\mathbb Z$-代数の存在、
(ii) $D^b(\operatorname{mod}B_R)$ 内で適切な相対 Abel heart を同定すること、
(iii) その heart を忘れた際の $\Pic$ 作用を関係式の非線形座標で表すこと、
は別の問題として残る。
特に本節の完備性を、GIT 粗商から無標識 Abel 圏への大域射の存在と
混同してはならない。
\end{remark}

\section{稠密な AS 型 locus と境界で不変な数値的 $E_6$}
\label{sec:generic-as-numerical}
本節は AOU 関係式から無標識 Abel 圏を構成できる範囲を二つの方向から限定する。
第一に、既知の楕円データによる構成は\emph{稠密開部分}で実際の AS 型 heart を与える。
第二に、Euler 形式と数値的 Serre 作用は全関係式空間にわたって一定であり、
これらだけから三次曲面型 heart を認識できないことを示す。

\subsection{楕円 decuple による実際の heart が存在する開部分}
\begin{proposition}[既知の双有理構成から得られる稠密 AS locus]
\label{prop:generic-as-locus}
$\mathbb C$ 上の AOU 関係式トーラス $M_Q^\circ$ には、空でない
Zariski 開部分 $U_{\rm ell}$ が存在して、各 $r\in U_{\rm ell}(\mathbb C)$ に対し、
滑らかな楕円曲線上の許容 decuple $\bigl(Y_r,(L_{v,r})_{v\in Q_0}\bigr)$ と
それから構成される AS-regular $\{0,1,2\}\times\mathbb Z$-代数
$A_r$ を選ぶことができる。このとき
\[
 \mathcal C_r^{\rm AS}:=\qgr A_r,\qquad
 D^b(\mathcal C_r^{\rm AS})\simeq\operatorname{Perf}(B_r),
\]
であり、右辺の九対象の有限 End 代数は $B_r$ である。
この主張は AOU の楕円 decuple との双有理対応と、その AS-regular 構成からの
\emph{既知結果の帰結}である。
特に $U_{\rm ell}$ の幾何学的点で三次曲面型 heart の候補が
存在することは、単なる形式 Artin 近傍の議論によらずに分かる。
ただし $\mathcal C_r^{\rm AS}$ の無標識同型類の一意性、
全 GIT 境界への延長、さらに定義\ref{def:weak-cubic}の追加的な
rank-one torsion-free 条件までは主張しない。
\end{proposition}
\begin{proof}
AOU\cite[Proposition~7.3, Theorem~1.1]{AOU} により許容楕円 decuple の
既約八次元モジュライ $M_{\rm ell}$ は、関係式の GIT 商 $M_Q^\chi$ と
双有理である。従って双方の空でない開部分で双有理写像と逆写像は正則同型になる。
この開部分を AOU の稠密トーラス $M_Q^\circ$ と交わらせ、
必要ならさらに縮小して $U_{\rm ell}$ とする。
AOU\cite[Lemma~7.4 and the paragraph following it]{AOU} が
岡川--植田の AS-regular 構成を適用することにより、
許容楕円 decuple は AS-regular $\{0,1,2\}\times\mathbb Z$-代数を与える。
この代数の標準九対象の End 代数は decuple の End 代数と一致し、
関係式点 $r$ に対応する $B_r$ となる。
AOU\cite[\S\S4.6, 7]{AOU} が構成する標準九対象は完全強例外列であり、
その有限 End 代数を用いる標準の tilting--Morita 同値から指定した導来同値を得る
（\cite[Definition~1.1]{OU2025} および AOU\cite[Theorem~9.4]{AOU} の具体的比較も参照）。
この論証は幾何学的点ごとの存在を示すものであり、
相対圏の非自明な降下や、異なる decuple から得る heart の一致を含意しない。
\end{proof}

\subsection{全関係式点に共通する Serre 形式と $E_6$}
$J$ をすべての成分が $1$ である $3\times3$ 行列、$I$ を単位行列とする。
右射影 $P_v=e_vB_r$ のクラスを層 $U,W,Z$ の順に並べると、
Euler 形式の行列（転置の規約に応じて入れ替わる）は
\begin{equation}\label{eq:aou-constant-euler}
 C=\begin{pmatrix}I&0&0\\J&I&0\\2J&J&I\end{pmatrix}
 \quad\text{または}\quad C^{\mathsf T}
\end{equation}
となる。これは各始点・終点間の道の次元だけから定まる。
数値的 Serre 作用を $S_r$ とすると Euler 双対性
$\chi(x,y)=\chi(y,S_rx)$ により、規約を一つ固定すれば
$S_r=C^{-1}C^{\mathsf T}$ である。

\begin{theorem}[普遍的な数値的 Serre 作用と $E_6$]
\label{thm:aou-numerical-E6}
任意の $r\in R_Q(\mathbb C)$ について
$K_0(\operatorname{Perf}B_r)$ の Euler 形式と Serre 作用は、
九つの標準射影加群の基底を用いて同じ整数行列で記述される。
特に $N_r:=S_r-1$ と置くと
\begin{equation}\label{eq:aou-unipotent-serre}
 \det(t-S_r)=(t-1)^9,\qquad
 \operatorname{rank}N_r=2,\qquad
 \operatorname{rank}N_r^2=1,\qquad N_r^3=0.
\end{equation}
さらに $\Lambda_r:=\ker N_r$ 上の Euler 形式は対称で、
その根基 $\operatorname{rad}(\chi|_{\Lambda_r})$ は階数一の原始部分格子である。
従って
\begin{equation}\label{eq:aou-numerical-root}
 \Lambda_r/\operatorname{rad}(\chi|_{\Lambda_r})
 \simeq Q(E_6)
\end{equation}
が\emph{正定値}の Euler 形式付き整数格子として成立する。
ここで右辺の形式は $E_6$ 根格子の負 Cartan 形式の符号を変えたもの、
すなわち判別式 $3$ の正定値形式である。
この数値的 $E_6$ は、AS 型 heart が未認識の境界点でも存在する。
\end{theorem}
\begin{proof}
定理\ref{thm:universal-directed-family}の階数計算から、
射影加群間の Hom は、同じ層なら恒等写像のみ、隣接層なら一組につき
一次元、最初から最後の層へは二次元である。
$B_r$ は有限大域次元なので射影間の高次 Ext は零であり、
式\eqref{eq:aou-constant-euler}を得る。各ファイバーで基底を
対応させると Euler 形式の整数行列は同一であり、
Serre 双対性から $S_r$ も同一である。
$J^2=3J$ を使い $C^{-1}C^{\mathsf T}$ を直接計算すれば
式\eqref{eq:aou-unipotent-serre}が従う。
再現可能な整数計算は
\path{computations/verify_aou_relation_tangent.py} にも収録した。

最後の格子の同定は、滑らかな可換三次曲面 $X$ の AOU 点 $r_0$ で
証明すればよい。AOU の完全強例外列は
$K_0(\operatorname{Perf}B_{r_0})\simeq K_0(X)$ を Euler 形式・Serre 作用とともに同一視する。
$X$ 上では $S_{r_0}$ は $(-)\otimes\omega_X[2]$ だから、
K 群上では $(-)\otimes\omega_X$ となる。
$\operatorname{ch}(x)=(a,D,z)$ と記すと
\[
 \operatorname{ch}((S_{r_0}-1)x)
   =(0,aK_X,D\cdot K_X+aK_X^2/2).
\]
$K_X$ は非零であり $K_0(X)$ は torsion-free なので、
$(S_{r_0}-1)x=0$ であるための必要十分条件は
$a=0$ かつ $D\cdot K_X=0$ である。
$\operatorname{rk}=0$ の二つのクラス $x,y$ に対して
Riemann--Roch から $\chi(x,y)=-c_1(x)\cdot c_1(y)$ となる。
従ってこの制限形式の根基は $c_1=0$ のクラス、すなわち
$\mathbb Z[\mathcal O_p]$ であり、これは原始である。
その商は $K_X^\perp\subset\operatorname{Pic}(X)$ に
負の交点形式の符号を変えて同型となる。
三次曲面が $\mathbb P^2$ の六点 blow-up であることから
$K_X^\perp$ は $E_6$ 根格子である
（命題\ref{prop:picard-root-index}）。
この同定を全 $r$ へ輸送すればよい。
\end{proof}

\begin{proposition}[矢の基底変換と Weyl 標識変更の相違]
\label{prop:aou-numerical-trivial-torus}
関係式 stack $[R_Q/G_Q]$ の矢の基底変換群 $G_Q$ は、
各 $B_r$ の九つの原始冪等元およびそれに対応する射影加群のクラスを固定する。
従って $G_Q$ は、全関係式空間にわたる数値的 $K_0$ と
式\eqref{eq:aou-numerical-root}の $E_6$ 格子に\emph{自明}に作用する。
特に可換 locus の古典的な非自明の $W(E_6)$ 標識変更は、
この矢の基底変換作用と同一視できない。
\end{proposition}
\begin{proof}
$G_Q$ は各矢の基底をスカラー倍するが、頂点冪等元 $e_v$ は固定する。
この代数同型による射影加群の輸送は $e_vB_r$ を
同じ添字 $v$ の射影加群へ送る。
従って $K_0$ の全標準基底を固定し、
そこから作った数値的 Serre 作用の固定部分格子と
その根基による商にも自明に作用する。
一方、古典的 $W(E_6)$ は根格子への忠実な反射作用を持つので、
非自明な標識変更を矢のスカラー倍だけで表すことはできない。
\end{proof}

\begin{corollary}[数値的不変量が区別できない境界]
\label{cor:aou-numerical-not-enough}
すべての関係式の係数が非零の点 $r_{\rm tor}$ と、
すべての九つの関係式が単項式である点 $r_{\rm mon}$ を取る。
この二つの smooth proper dg 圏は、Euler 形式・数値的 Serre 作用・
式\eqref{eq:aou-numerical-root}の $E_6$ 格子について同じ不変量を持つが、
\[
 \dim HH^1(B_{r_{\rm tor}})=0,\qquad
 \dim HH^1(B_{r_{\rm mon}})=10
\]
である。従って $\operatorname{Perf}(B_{r_{\rm tor}})$ と
$\operatorname{Perf}(B_{r_{\rm mon}})$ は導来 Morita 同値ではない。
特に Euler 形式と数値的 $E_6$ 標識だけでは、
関係式空間の異なる導来圏、ひいては三次曲面型 Abel heart の
認識問題を解決できない。
\end{corollary}
\begin{proof}
前定理の数値的不変量の一定性と
定理\ref{thm:aou-hh-support}の Hochschild 次元公式を組み合わせる。
Hochschild cohomology は導来 Morita 同値で保存されるので、
$HH^1$ の次元が異なる二圏は Morita 同値にならない。
これは単項式点に\emph{いかなる}三次曲面型 heart も存在しないと
主張するものではない。
\end{proof}

\begin{proposition}[Hochschild 次元が跳ばない境界点]
\label{prop:aou-boundary-hh-minimal}
九関係式のうち、$(u_0,z_0)$ の関係式だけを $[1:0:0]$ へ特殊化し、
他の八関係式はすべての係数が非零であるとする。
この点は稠密トーラスの外にあるが
\[
 \rho(r)=10,\qquad
 (\dim HH^1(B_r),\dim HH^2(B_r))=(0,8)
\]
が成り立つ。従って $HH^1,HH^2$ の次元が一般値に等しいことすら、
全関係式係数の非零性を特徴付けない。
\end{proposition}
\begin{proof}
定理\ref{thm:aou-hh-support}の支持行列を用いる。
中間層の各 $m=1,2$ に対して
$x_i=a_{im}-a_{i0}$、$y_j=b_{mj}-b_{0j}$ と書くと、
残り八つの関係式は $x_i+y_j=0$ を $(i,j)\ne(0,0)$ の
すべてで要求する。二部グラフ $K_{3,3}$ から辺 $(0,0)$ だけを
除いたグラフは連結であるため、これらから $x_i=-y_j$ が
\emph{すべて}の $(i,j)$ について従う。
よって欠けた $(0,0)$ の二つの関係式は他の十六行の線形結合であり、
支持行列の階数は全係数非零のときと同じ $10$ である。
最後の式は Hochschild 次元公式より従う。
\end{proof}

\begin{remark}[今後の標識忘却への含意]
\label{rem:numerical-unmark}
数値的 Serre 作用と $E_6$ は全関係式族に共通して存在するが、
幾何学的に許容される helix、三次曲面型 heart、
線束に類する構造対象およびそれらの exact 自己同値作用は、
この格子だけから復元できない。
従って AOU の\emph{数値的標識}から\emph{幾何学的標識}への
認識が、無標識 Abel 圏への大域降下に先立つ独立の課題である。
稠密な AS 型 locus の存在（命題\ref{prop:generic-as-locus}）と
全境界における数値的 $E_6$ の一定性（定理\ref{thm:aou-numerical-E6}）を
混同してはならない。
\end{remark}

\section{可換 AOU 点の形式変形と集合論的な標識変更}
\label{sec:formal-abelian-recognition}
本節は固定された可換反標準 heart の変形理論の補助結果を記す。
得られる形式 Abel 変形が反標準 heart の許容条件を満たす範囲を、
抽象的な Abel 変形の全体と区別する。
鍵になるのはDe Deken--Lowen\cite[Theorems 4.5, 5.13, 5.15]{DDL}であり、
本稿の有限代数変形の完備性（定理\ref{thm:aou-relation-formal-completeness}）
と組み合わせる。導来圏の任意のheartを自動的に認識するのではなく、
特定の可換三次曲面から始まる\emph{Artin形式変形}を扱う。

\subsection{標準helixとAbel変形の比較}
$X$を滑らかな複素三次曲面とし、
$E_0,\ldots,E_8$をAOUの標準九対象line-bundle threadとする。
$\mathscr H=(E_i)_{i\in\mathbb Z}$は、このthreadが生成する
幾何学的helixである。以下、文献でいう
\emph{$\mathbb Z$-generating geometric helix}の条件、すなわち
$(E_i)$が$\operatorname{Qcoh}X$を$\mathbb Z$-生成し、
$i\le j$について$\operatorname{Ext}^{1,2}(E_i,E_j)=0$を満たすことを用いる。
標準三ブロックhelixの幾何学的性質は
Karpov--Nogin\cite{KN}およびAOU\cite[\S4]{AOU}の既知の入力である。
$-K_X$の豊富性と$E_{i+9}\cong E_i\otimes\mathcal O_X(-K_X)$により、
このhelixの適切な反標準方向の部分列が生成列となる。

\begin{theorem}[AOU形式関係式からのAbel変形の再構成]
\label{thm:aou-formal-abelian-reconstruction}
$X,\mathscr H$を上のとおりとし、$r\in R_Q$をそのAOU関係式点、
$B=\operatorname{End}_X(\bigoplus_{i=0}^{8}E_i)$とする。
$R$を剰余体$\mathbb C$の局所Artin $\mathbb C$-代数とする。
$r$に特殊化する任意のAOU関係式の$R$-点$r_R$について、
付随する平坦な有限代数$B_R$から、
$\mathcal C_R\otimes_R\mathbb C\simeq\operatorname{Qcoh}X$という
特殊ファイバー同定を伴う平坦なGrothendieck Abel変形
$\mathcal C_R$が構成できる。
標準threadの変形を保持したAbel変形の同型類集合と、
九つの原始冪等元を標識した$B$の平坦な有限代数変形の同型類集合は
自然に全単射である。これらの同値は$R$の小拡大と底変更に両立する。
さらに$\mathcal C_R$は局所noetherianであり、
そのnoetherian対象の部分圏は$\operatorname{coh}X$からの
Abel変形を与える。
\end{theorem}
\begin{proof}
標準helixから作る$\mathbb Z$-代数を$\mathfrak a$、
九対象threadへの制限を$\mathfrak a_{[0,8]}$と書く。
後者は$B$の原始冪等元を標識した線形カテゴリーである。
De Deken--Lowen\cite[Theorem 5.15]{DDL}は、上で述べた生成性、
幾何学的helixおよびExt消滅のもとで
\[
 \operatorname{Def}_{\rm ab}(\operatorname{Qcoh}X)
 \simeq\operatorname{Def}_{\rm lin}(\mathfrak a)
 \simeq\operatorname{Def}_{\rm lin}(\mathfrak a_{[0,8]})
\]
を与える。第二の変形問題は、原始冪等元を標識した
有限代数$B$の平坦な変形問題と等しい。
定理\ref{thm:aou-relation-formal-completeness}により
この変形は九つの関係式の平坦な変形により尽くされる。
上記の同値の逆を適用すれば$\mathcal C_R$を得る。
この議論はAbel変形の同値を使うため、
単なる$D^b(\operatorname{mod}B_R)$の標準heartを
$\operatorname{coh}X$と取り違えてはいない。

局所noetherian性も確認する。De Deken--Lowenの再構成では
$\mathbb Z$-生成列$(E_i)$が$R$-flatな生成対象$(E_{i,R})$へ持ち上がる。
$R$の極大イデアルを$\mathfrak m$とすると
$E_{i,R}$の有限$\mathfrak m$-adic filtrationの各商は
$E_i$の有限直和である。$E_i$はnoetherianであるから
$E_{i,R}$もnoetherianである。
従って$\mathcal C_R$はnoetherian生成対象を持つ
局所noetherian Grothendieck圏であり、noetherian対象の部分圏は
核・余核・有限拡大で閉じる。特殊ファイバーのnoetherian対象は
$\operatorname{coh}X$となる。
\end{proof}

\begin{remark}[復元の適用範囲]
\label{rem:aou-formal-heart-scope}
この定理は\emph{滑らかな可換点を固定した任意の局所Artin底}での
Abel変形の存在を与える。$R_Q$の全閉点で三次曲面型のheartが
代数的族として存在するという主張ではない。
またGrothendieck Abel変形の存在だけから
$\mathsf S_R^{-1}[2]$の$t$-exact性、rank-one構造対象の公理、
Artin--Zhang豊富性が形式変形の全域で自動的に成立すると
断言しない。これらを要する本稿の弱標識付きモジュライへの
因子化には、さらに検証を要する。
\end{remark}

\subsection{形式八次元チャートと標識変更}
$\mathcal D_X$を$\operatorname{Qcoh}X$の平坦なAbel変形関手で、
特殊ファイバーとの同定を保持し、
関手の自然同型を同一視したものとする。

\begin{theorem}[AOUの形式八次元チャート]
\label{thm:aou-formal-chart}
上の$X,r$について、$r$はAOUの全係数非零のトーラス
$M_Q^\circ$の点であるとする。
標準helixの選択に依存する形式変形関手の同型
\[
 \widehat{M_Q^\circ}_{[r]}\simeq \mathcal D_X
\]
が存在する。従って特殊ファイバーを同定したAbel変形の
同型類は、$\mathbb C$上の八次元形式円板で局所的に記述される。
その接空間は
\[
 T_{[r]}M_Q^\circ\simeq HH^2(B)\simeq HH^2(X)
 =H^1(X,T_X)\oplus H^0(X,-K_X)
\]
である。形式チャートの同型は標準helixと特殊化の同定に依存し、
無標識圏の大域的な正則パラメータ表示ではない。
\end{theorem}
\begin{proof}
定理\ref{thm:aou-formal-abelian-reconstruction}により
Abel変形は有限代数変形と同値であり、
定理\ref{thm:aou-relation-formal-completeness}により
九関係式の変形を矢の基底変更で割ったものと同値である。
全係数非零のトーラスでは、矢の基底変更群の
関係式への作用の核は頂点スカラーによる八次元トーラスである。
有効商トーラスの作用は自由である。
実際、作用行列のSmith標準形が
$\operatorname{diag}(1^{\times10},0^{\times8})$となることは、
命題\ref{prop:aou-kodaira-spencer-iso}の格子計算からも分かる。
従って関係式変形の同型類関手は
滑らかなトーラス商$M_Q^\circ$の点$[r]$における形式完成となる。
接空間については命題\ref{prop:aou-kodaira-spencer-iso}を適用する。
\end{proof}

\begin{proposition}[形式Picard作用と一次shear]
\label{prop:formal-picard-action}
$X$のexact自己同値群
$G_X=\operatorname{Pic}(X)\rtimes\operatorname{Aut}(X)$は、
特殊ファイバーの同定を変更することにより
$\mathcal D_X$に作用する。前定理の形式チャートへ輸送すると、
$\widehat{M_Q^\circ}_{[r]}$上に
一般には非線形な\emph{形式的}作用が定まる。
その$L\in\operatorname{Pic}(X)$による接空間作用は、
方向の規約による符号を除いて
\[
 (\xi,\pi)\longmapsto
 (\xi+H^1(\pi^\sharp)c_1(L),\pi)
\]
である。これは命題\ref{prop:picard-first-order}のshearに等しい。
この形式作用からAOUの大域GITスキームへのPicard群の
正則作用は導かれない。
\end{proposition}
\begin{proof}
特殊ファイバーとの同定を$\phi$から$F\circ\phi$へ変更する操作は
各$F\in\operatorname{Aut}_{\rm ex}(\operatorname{Qcoh}X)$について
Abel変形関手の自己同型を与え、合成に関して群作用となる。
可換滑らかな射影多様体のAbel圏の自己同値は
自己同型の引戻しと線束捩れで記述され、
$G_X=\operatorname{Pic}(X)\rtimes\operatorname{Aut}(X)$である。
定理\ref{thm:aou-formal-chart}によって作用を形式完成上に輸送する。
一次項は線束Morita変換のČech計算
（命題\ref{prop:picard-first-order}）により与えられる。
これはToda\cite{Toda}のHochschild変形に対する
Fourier--Mukai変換とも整合する。
\end{proof}

\begin{corollary}[無標識形式変形の同値類集合]
\label{cor:aou-formal-unmarked-quotient}
特殊ファイバーとの同定を忘れた Abel 変形の exact 同値類集合は、
同定付き変形の同型類集合 $\mathcal D_X$ の
$G_X=\operatorname{Pic}(X)\rtimes\operatorname{Aut}(X)$ による軌道集合である。
この主張は集合値の形式変形問題についてのものである。
反標準 heart と弱標識の条件を必要とする場合は、その条件を満たす部分に制限する。
\end{corollary}
\begin{proof}
基礎圏の二つの同定の差は、その特殊ファイバーの exact 自己同値である。
逆に同定の変更は、同定を忘れた圏の同値類を変えない。
従って同値関係はちょうどこの軌道関係である。
\end{proof}

\begin{remark}[共通線束捩れと形式チャートの非標準性]
\label{rem:formal-common-twist}
同じ$X$でthreadを$(E_i\otimes L)$へ取り替えても
$\operatorname{End}(\bigoplus_iE_i\otimes L)$は$B$と同一視され、
AOU関係式点は変わらない（命題\ref{prop:common-tensor}）。
しかしDe Deken--Lowenの逆構成を特殊ファイバー同定付きで比較すると、
二つの形式チャートは線束捩れによる再標識で移り合う。
その一次遷移は命題\ref{prop:formal-picard-action}のshearである。
従って関係式点から$\mathcal O$の絶対的な選択が自動的に
得られるわけではない。まさにこの差が無標識Abel圏を考える動機となる。
\end{remark}


\section{再現性と検証の範囲}
計算コードは \path{computations/verify_marking_lattice.py}、
\path{computations/verify_aou_relation_tangent.py}、
\path{computations/verify_line_bundle_numerics.py}、
\path{computations/verify_weyl_frames.py}、
静的参照検査は \path{scripts/verify_unmark_bundle.py} である。
計算は数値恒等式を検算するもので、線束の存在やheartの同定の証明ではない。
反標準 Serre 消滅・階数認識・root 型障害・格子標識復元の証明を点検した。
数値上の有限探索では27本の $(-1)$ 類の網羅性、72個の六組、
$51840$ 個の順序付き標識、単純反射の Coxeter 関係を検算する。
これらは三対象交換の圏内での実現を証明しない。
保持した既存証明すべての独立再査読を実施したという意味でもない。
検証状況は \path{research/unmark_validation.json} に記す。
\begin{thebibliography}{99}
\bibitem{AZ} M. Artin, J. J. Zhang,
\emph{Noncommutative projective schemes},
Adv. Math. 109 (1994), 228--287,
\href{https://doi.org/10.1006/aima.1994.1087}{doi:10.1006/aima.1994.1087}.
参照: Theorem 4.5（ample pairからの再構成）。
\bibitem{AOU} T. Abdelgadir, S. Okawa, K. Ueda,
\emph{Compact moduli of marked noncommutative cubic surfaces},
\href{https://arxiv.org/abs/2404.00175}{arXiv:2404.00175}, v1.
参照: Introduction, \S\S4.6, 5.1--5.4, Theorems 1.1, 1.2, 9.4, Lemma 9.2.
\bibitem{Lowen} W. Lowen,
\emph{Obstruction theory for objects in abelian and derived categories},
Comm. Algebra 33 (2005), 3195--3223,
\href{https://arxiv.org/abs/math/0407019}{arXiv:math/0407019}, Theorem A.
\bibitem{KN} B. V. Karpov, D. Yu. Nogin,
\emph{Three-block exceptional collections over Del Pezzo surfaces},
Izv. Math. 62 (1998), 429--463,
\href{https://doi.org/10.1070/IM1998v062n03ABEH000205}{doi:10.1070/IM1998v062n03ABEH000205}.
参照: \S1.3（可換例外層の数値的一意性）、三ブロック列の分類。
\bibitem{DDL} O. De Deken, W. Lowen,
\emph{Abelian and derived deformations in the presence of $\mathbb Z$-generating geometric helices},
J. Noncommut. Geom. 5 (2011), 477--505,
\href{https://doi.org/10.4171/JNCG/83}{doi:10.4171/JNCG/83},
Theorems 4.5, 5.13, 5.15 and Definition 5.12.
\bibitem{OU2025} S. Okawa, K. Ueda,
\emph{Marked Artin--Schelter surfaces of del Pezzo types},
\href{https://arxiv.org/abs/2508.04986}{arXiv:2508.04986v1} (2025),
Definition 1.1 and Theorem 6.4.
\bibitem{DGK} I. Dolgachev, B. van Geemen, S. Kond\=o,
\emph{A complex ball uniformization of the moduli space of cubic surfaces
via periods of K3 surfaces},
J. reine angew. Math. 588 (2005), 99--148,
\href{https://doi.org/10.1515/crll.2005.2005.588.99}{doi:10.1515/crll.2005.2005.588.99};
\S\S 2.5--2.6.
\bibitem{Gabriel} P. Gabriel, \emph{Des catégories abéliennes},
Bull. Soc. Math. France 90 (1962), 323--448,
\href{https://doi.org/10.24033/bsmf.1583}{doi:10.24033/bsmf.1583}.
\bibitem{BO} A. Bondal, D. Orlov,
\emph{Reconstruction of a variety from the derived category and groups of autoequivalences},
Compositio Math. 125 (2001), 327--344;
\href{https://arxiv.org/abs/alg-geom/9712029}{arXiv:alg-geom/9712029}.
\bibitem{Toda} Y. Toda, \emph{Deformations and Fourier--Mukai transforms},
J. Differential Geom. 81 (2009), 197--224;
\href{https://arxiv.org/abs/math/0502571}{arXiv:math/0502571}.
\bibitem{OWR} K. Ueda (joint work with S. Okawa),
\emph{Affine Weyl group actions on moduli spaces of marked noncommutative del Pezzo surfaces},
Oberwolfach Reports 44/2024, pp. 2557--2558,
\href{https://doi.org/10.4171/OWR/2024/44}{doi:10.4171/OWR/2024/44}, Conjectures 1--3.
\bibitem{PolishchukProj} A. Polishchuk,
\emph{Noncommutative Proj and coherent algebras},
Math. Res. Lett. 12 (2005), 63--74,
\href{https://arxiv.org/abs/math/0212182}{arXiv:math/0212182}.
参照: \S1、Lemma 2.1、Theorem 2.4。coherent と noetherian の区別を用いる。
\bibitem{HR} N. Hoffmann, F. Reede,
\emph{Torsion-free rank one sheaves over del Pezzo orders},
J. Algebra 493 (2018), 251--266,
\href{https://doi.org/10.1016/j.jalgebra.2017.08.029}{doi:10.1016/j.jalgebra.2017.08.029};
\href{https://arxiv.org/abs/1603.03324}{arXiv:1603.03324v1}.
参照: Lemma 3.9（同じ第一 Chern 類での第二 Ext 消滅とその証明）。
\bibitem{ST} P. Seidel, R. Thomas,
\emph{Braid group actions on derived categories of coherent sheaves},
\href{https://arxiv.org/abs/math/0001043}{arXiv:math/0001043v2} (2000).
参照: Definitions 2.7, 2.9、Proposition 2.10。評価複体の既知背景として引用する。
\end{thebibliography}
\end{document}
